
\documentclass[conference]{IEEEtran}
\IEEEoverridecommandlockouts
\usepackage{amsmath,amssymb,bm}
\usepackage{booktabs}
\usepackage{graphicx}
\usepackage{subfig}
\usepackage{multirow}
\usepackage{array}
\usepackage{siunitx}
\usepackage{hyperref}
\usepackage{enumitem}
\usepackage{listings}
\usepackage{xcolor}
\hypersetup{colorlinks=true,linkcolor=blue,citecolor=blue,urlcolor=blue}

\title{AMASA-Complex v2: Full-Length Technical Report on Theory, Architecture, Safety, and Empirical Analysis}
\author{\IEEEauthorblockN{Codex}\\
\IEEEauthorblockA{Workspace: \texttt{/Users/tahamajs/Documents/uni/DRL/projects/amasa\_clean}}}

\begin{document}
\maketitle

\begin{abstract}
This full-length report documents AMASA-Complex v2 as a complete educational and engineering platform for safety-critical reinforcement learning in autonomous suturing. Beyond presenting implementation details, the report explains the theoretical motivations behind each architectural decision, the operational failure modes encountered in practical training, and the reproducible evaluation process required for credible reward-safety claims. The project combines four learning algorithms (CQL, IQL, SAC-Lagrangian, PPO-Lagrangian), a hybrid guard stack (PID dual regulation, risk critic gating, and decision-tree shielding), and three scenario families for robustness stress testing. The resulting system is configuration-driven, CPU-first, and backward-compatible with legacy command-line workflows.

The report intentionally emphasizes methodological rigor. In this domain, performance should not be summarized by a single reward number; instead, it should be interpreted through Pareto behavior, threshold tiers, and telemetry traces that expose transient risk and control instability. We therefore provide benchmark matrix outcomes, expanded PID sweep analysis, and safety timeline interpretation, then link these observations back to concrete algorithmic and systems-level causes.

Finally, this document serves as a course-scale reference artifact. It is structured to support teaching, replication, and extension: mathematical background, design rationale, implementation mapping, validation evidence, and future ablation plans are all presented in one place.
\end{abstract}

\section{Introduction}
Autonomous surgical assistance is an ideal capstone domain for advanced reinforcement learning because reward maximization, safety guarantees, and robustness constraints all matter simultaneously. A model that achieves high reward with unstable safety behavior is not deployable. A model that never violates constraints but cannot complete tasks is also not useful. The central challenge is to push the reward-safety Pareto frontier, not to optimize one axis in isolation.

AMASA-Complex v2 is designed around this principle. It upgrades a lightweight benchmark into a system where students and researchers can analyze the full offline-to-online pipeline under controlled scenario perturbations. The project also addresses practical software concerns that often dominate real failures: config precedence mistakes, inconsistent checkpoint semantics, missing telemetry, and non-reproducible evaluation scripts.

\section{Theory Recap and Design Constraints}
\subsection{From MDP to CMDP}
Let $\mathcal{M}=(\mathcal{S},\mathcal{A},P,R,\gamma)$ denote the task MDP. Safety-critical control adds a cost channel with limit $d$:
\[
\max_{\pi} J_R(\pi)\;\text{s.t.}\;J_C(\pi)\le d.
\]
The Lagrangian objective is:
\[
\mathcal{L}(\pi,\lambda)=J_R(\pi)-\lambda(J_C(\pi)-d),\;\lambda\ge0.
\]
In practice, this objective is only as stable as the dual update rule and the quality of cost estimation.

\subsection{Offline Shift and Conservatism}
Offline datasets provide cheap initialization but impose support constraints. Both CQL and IQL in this project are used to produce high-quality starter policies while minimizing unsafe optimism under OOD actions. This is critical before entering costly online fine-tuning.

\subsection{Safety as Closed-Loop Control}
PID regulation of $\lambda$ is used to stabilize cost tracking:
\[
\lambda_{t+1}=\mathrm{clip}(K_p e_t + K_i\sum_{\tau\le t} e_\tau + K_d(e_t-e_{t-1}),0,\lambda_{\max}).
\]
This control view makes gain sweeps a first-class experimental procedure rather than an afterthought.

\section{Benchmark Definition and System Scope}
The environment keeps the 23-dimensional observation and 7-dimensional action interfaces while introducing scenario controls for force noise, dynamics shifts, observation corruption, and delayed observations. Three scenario families are used: nominal, perturbed, adversarial.

The implementation adds modular packages for core config/registry, environments, offline and online algorithms, safety stack, and benchmark aggregation. Legacy scripts remain operational with optional modern flags for config-driven execution.

\section{Comprehensive Concept Modules}
The following modules provide detailed conceptual explanations and engineering interpretation. Each module is written as a mini-lesson with actionable diagnostics.

\subsubsection{Module 1: Markov Modeling in Surgical Tasks (Foundational View)}
\textbf{Principle.} A stable MDP formalization is the first quality gate. In AMASA-Complex v2 this appears as a constraint on both algorithm design and experiment process. The key point is that no isolated formula is sufficient; correctness emerges from interaction between objective, data regime, guard behavior, and reproducibility checks.

\textbf{Failure mode.} A common practical breakdown is state aliasing across camera artifacts. This usually manifests as deceptive short-run gains followed by cross-seed instability. In safety-critical tasks, this instability must be interpreted as a primary failure signal even when average reward appears competitive.

\textbf{Mitigation.} The project mitigation is to explicitly separate physical state proxies from visual nuisance channels. This mitigation is encoded in configuration and code so that it can be repeated consistently across seeds and scenarios. It is also paired with evaluation constraints to avoid post-hoc metric selection bias.

\textbf{Depth guidance (foundational).} At the foundational level, students should be able to restate this concept formally, map it to concrete code paths, and identify the minimum telemetry needed to validate implementation correctness.

\textbf{Operational checklist.} Before accepting any run result for this module, verify: (i) metric definitions are unchanged across runs, (ii) scenario and seed lists are fixed, (iii) checkpoint metadata is parseable by automated evaluators, (iv) run logs include safety channels, and (v) threshold tier classification is reported together with raw reward and cost.

\begin{itemize}[leftmargin=*]
  \item Principle statement: A stable MDP formalization is the first quality gate.
  \item Primary challenge: state aliasing across camera artifacts.
  \item Mitigation path: explicitly separate physical state proxies from visual nuisance channels.
  \item Verification channels: reward, cost, lambda, risk, block rate.
  \item Success criterion: positive frontier trend with stable safety traces.
  \item Documentation rule: report both mean and seed-level outcomes.
\end{itemize}


\subsubsection{Module 2: Markov Modeling in Surgical Tasks (Advanced Research View)}
\textbf{Principle.} A stable MDP formalization is the first quality gate. In AMASA-Complex v2 this appears as a constraint on both algorithm design and experiment process. The key point is that no isolated formula is sufficient; correctness emerges from interaction between objective, data regime, guard behavior, and reproducibility checks.

\textbf{Failure mode.} A common practical breakdown is state aliasing across camera artifacts. This usually manifests as deceptive short-run gains followed by cross-seed instability. In safety-critical tasks, this instability must be interpreted as a primary failure signal even when average reward appears competitive.

\textbf{Mitigation.} The project mitigation is to explicitly separate physical state proxies from visual nuisance channels. This mitigation is encoded in configuration and code so that it can be repeated consistently across seeds and scenarios. It is also paired with evaluation constraints to avoid post-hoc metric selection bias.

\textbf{Depth guidance (advanced).} At the advanced level, the same concept should be treated as a research variable: students should formulate competing hypotheses, design controlled ablations, and reason about how uncertainty in one subsystem propagates into frontier movement.

\textbf{Operational checklist.} Before accepting any run result for this module, verify: (i) metric definitions are unchanged across runs, (ii) scenario and seed lists are fixed, (iii) checkpoint metadata is parseable by automated evaluators, (iv) run logs include safety channels, and (v) threshold tier classification is reported together with raw reward and cost.

\begin{itemize}[leftmargin=*]
  \item Principle statement: A stable MDP formalization is the first quality gate.
  \item Primary challenge: state aliasing across camera artifacts.
  \item Mitigation path: explicitly separate physical state proxies from visual nuisance channels.
  \item Verification channels: reward, cost, lambda, risk, block rate.
  \item Success criterion: positive frontier trend with stable safety traces.
  \item Documentation rule: report both mean and seed-level outcomes.
\end{itemize}


\subsubsection{Module 3: Constrained Optimization in CMDPs (Foundational View)}
\textbf{Principle.} Safety cost is not a soft preference but a hard reliability channel. In AMASA-Complex v2 this appears as a constraint on both algorithm design and experiment process. The key point is that no isolated formula is sufficient; correctness emerges from interaction between objective, data regime, guard behavior, and reproducibility checks.

\textbf{Failure mode.} A common practical breakdown is dual variables can oscillate when costs are sparse. This usually manifests as deceptive short-run gains followed by cross-seed instability. In safety-critical tasks, this instability must be interpreted as a primary failure signal even when average reward appears competitive.

\textbf{Mitigation.} The project mitigation is to PID-based lambda control dampens oscillations with derivative action. This mitigation is encoded in configuration and code so that it can be repeated consistently across seeds and scenarios. It is also paired with evaluation constraints to avoid post-hoc metric selection bias.

\textbf{Depth guidance (foundational).} At the foundational level, students should be able to restate this concept formally, map it to concrete code paths, and identify the minimum telemetry needed to validate implementation correctness.

\textbf{Operational checklist.} Before accepting any run result for this module, verify: (i) metric definitions are unchanged across runs, (ii) scenario and seed lists are fixed, (iii) checkpoint metadata is parseable by automated evaluators, (iv) run logs include safety channels, and (v) threshold tier classification is reported together with raw reward and cost.

\begin{itemize}[leftmargin=*]
  \item Principle statement: Safety cost is not a soft preference but a hard reliability channel.
  \item Primary challenge: dual variables can oscillate when costs are sparse.
  \item Mitigation path: PID-based lambda control dampens oscillations with derivative action.
  \item Verification channels: reward, cost, lambda, risk, block rate.
  \item Success criterion: positive frontier trend with stable safety traces.
  \item Documentation rule: report both mean and seed-level outcomes.
\end{itemize}


\subsubsection{Module 4: Constrained Optimization in CMDPs (Advanced Research View)}
\textbf{Principle.} Safety cost is not a soft preference but a hard reliability channel. In AMASA-Complex v2 this appears as a constraint on both algorithm design and experiment process. The key point is that no isolated formula is sufficient; correctness emerges from interaction between objective, data regime, guard behavior, and reproducibility checks.

\textbf{Failure mode.} A common practical breakdown is dual variables can oscillate when costs are sparse. This usually manifests as deceptive short-run gains followed by cross-seed instability. In safety-critical tasks, this instability must be interpreted as a primary failure signal even when average reward appears competitive.

\textbf{Mitigation.} The project mitigation is to PID-based lambda control dampens oscillations with derivative action. This mitigation is encoded in configuration and code so that it can be repeated consistently across seeds and scenarios. It is also paired with evaluation constraints to avoid post-hoc metric selection bias.

\textbf{Depth guidance (advanced).} At the advanced level, the same concept should be treated as a research variable: students should formulate competing hypotheses, design controlled ablations, and reason about how uncertainty in one subsystem propagates into frontier movement.

\textbf{Operational checklist.} Before accepting any run result for this module, verify: (i) metric definitions are unchanged across runs, (ii) scenario and seed lists are fixed, (iii) checkpoint metadata is parseable by automated evaluators, (iv) run logs include safety channels, and (v) threshold tier classification is reported together with raw reward and cost.

\begin{itemize}[leftmargin=*]
  \item Principle statement: Safety cost is not a soft preference but a hard reliability channel.
  \item Primary challenge: dual variables can oscillate when costs are sparse.
  \item Mitigation path: PID-based lambda control dampens oscillations with derivative action.
  \item Verification channels: reward, cost, lambda, risk, block rate.
  \item Success criterion: positive frontier trend with stable safety traces.
  \item Documentation rule: report both mean and seed-level outcomes.
\end{itemize}


\subsubsection{Module 5: Offline Distribution Shift (Foundational View)}
\textbf{Principle.} policy evaluation must stay conservative outside the dataset support. In AMASA-Complex v2 this appears as a constraint on both algorithm design and experiment process. The key point is that no isolated formula is sufficient; correctness emerges from interaction between objective, data regime, guard behavior, and reproducibility checks.

\textbf{Failure mode.} A common practical breakdown is Q-functions overestimate unseen actions. This usually manifests as deceptive short-run gains followed by cross-seed instability. In safety-critical tasks, this instability must be interpreted as a primary failure signal even when average reward appears competitive.

\textbf{Mitigation.} The project mitigation is to conservative regularization and behavior-closeness penalties reduce OOD exploitation. This mitigation is encoded in configuration and code so that it can be repeated consistently across seeds and scenarios. It is also paired with evaluation constraints to avoid post-hoc metric selection bias.

\textbf{Depth guidance (foundational).} At the foundational level, students should be able to restate this concept formally, map it to concrete code paths, and identify the minimum telemetry needed to validate implementation correctness.

\textbf{Operational checklist.} Before accepting any run result for this module, verify: (i) metric definitions are unchanged across runs, (ii) scenario and seed lists are fixed, (iii) checkpoint metadata is parseable by automated evaluators, (iv) run logs include safety channels, and (v) threshold tier classification is reported together with raw reward and cost.

\begin{itemize}[leftmargin=*]
  \item Principle statement: policy evaluation must stay conservative outside the dataset support.
  \item Primary challenge: Q-functions overestimate unseen actions.
  \item Mitigation path: conservative regularization and behavior-closeness penalties reduce OOD exploitation.
  \item Verification channels: reward, cost, lambda, risk, block rate.
  \item Success criterion: positive frontier trend with stable safety traces.
  \item Documentation rule: report both mean and seed-level outcomes.
\end{itemize}


\subsubsection{Module 6: Offline Distribution Shift (Advanced Research View)}
\textbf{Principle.} policy evaluation must stay conservative outside the dataset support. In AMASA-Complex v2 this appears as a constraint on both algorithm design and experiment process. The key point is that no isolated formula is sufficient; correctness emerges from interaction between objective, data regime, guard behavior, and reproducibility checks.

\textbf{Failure mode.} A common practical breakdown is Q-functions overestimate unseen actions. This usually manifests as deceptive short-run gains followed by cross-seed instability. In safety-critical tasks, this instability must be interpreted as a primary failure signal even when average reward appears competitive.

\textbf{Mitigation.} The project mitigation is to conservative regularization and behavior-closeness penalties reduce OOD exploitation. This mitigation is encoded in configuration and code so that it can be repeated consistently across seeds and scenarios. It is also paired with evaluation constraints to avoid post-hoc metric selection bias.

\textbf{Depth guidance (advanced).} At the advanced level, the same concept should be treated as a research variable: students should formulate competing hypotheses, design controlled ablations, and reason about how uncertainty in one subsystem propagates into frontier movement.

\textbf{Operational checklist.} Before accepting any run result for this module, verify: (i) metric definitions are unchanged across runs, (ii) scenario and seed lists are fixed, (iii) checkpoint metadata is parseable by automated evaluators, (iv) run logs include safety channels, and (v) threshold tier classification is reported together with raw reward and cost.

\begin{itemize}[leftmargin=*]
  \item Principle statement: policy evaluation must stay conservative outside the dataset support.
  \item Primary challenge: Q-functions overestimate unseen actions.
  \item Mitigation path: conservative regularization and behavior-closeness penalties reduce OOD exploitation.
  \item Verification channels: reward, cost, lambda, risk, block rate.
  \item Success criterion: positive frontier trend with stable safety traces.
  \item Documentation rule: report both mean and seed-level outcomes.
\end{itemize}


\subsubsection{Module 7: Causal Feature Separation (Foundational View)}
\textbf{Principle.} intervention-invariant features improve robustness. In AMASA-Complex v2 this appears as a constraint on both algorithm design and experiment process. The key point is that no isolated formula is sufficient; correctness emerges from interaction between objective, data regime, guard behavior, and reproducibility checks.

\textbf{Failure mode.} A common practical breakdown is spurious visual correlations leak into value estimates. This usually manifests as deceptive short-run gains followed by cross-seed instability. In safety-critical tasks, this instability must be interpreted as a primary failure signal even when average reward appears competitive.

\textbf{Mitigation.} The project mitigation is to introduce data corruption stress tests and track degradation under lighting/camera shifts. This mitigation is encoded in configuration and code so that it can be repeated consistently across seeds and scenarios. It is also paired with evaluation constraints to avoid post-hoc metric selection bias.

\textbf{Depth guidance (foundational).} At the foundational level, students should be able to restate this concept formally, map it to concrete code paths, and identify the minimum telemetry needed to validate implementation correctness.

\textbf{Operational checklist.} Before accepting any run result for this module, verify: (i) metric definitions are unchanged across runs, (ii) scenario and seed lists are fixed, (iii) checkpoint metadata is parseable by automated evaluators, (iv) run logs include safety channels, and (v) threshold tier classification is reported together with raw reward and cost.

\begin{itemize}[leftmargin=*]
  \item Principle statement: intervention-invariant features improve robustness.
  \item Primary challenge: spurious visual correlations leak into value estimates.
  \item Mitigation path: introduce data corruption stress tests and track degradation under lighting/camera shifts.
  \item Verification channels: reward, cost, lambda, risk, block rate.
  \item Success criterion: positive frontier trend with stable safety traces.
  \item Documentation rule: report both mean and seed-level outcomes.
\end{itemize}


\subsubsection{Module 8: Causal Feature Separation (Advanced Research View)}
\textbf{Principle.} intervention-invariant features improve robustness. In AMASA-Complex v2 this appears as a constraint on both algorithm design and experiment process. The key point is that no isolated formula is sufficient; correctness emerges from interaction between objective, data regime, guard behavior, and reproducibility checks.

\textbf{Failure mode.} A common practical breakdown is spurious visual correlations leak into value estimates. This usually manifests as deceptive short-run gains followed by cross-seed instability. In safety-critical tasks, this instability must be interpreted as a primary failure signal even when average reward appears competitive.

\textbf{Mitigation.} The project mitigation is to introduce data corruption stress tests and track degradation under lighting/camera shifts. This mitigation is encoded in configuration and code so that it can be repeated consistently across seeds and scenarios. It is also paired with evaluation constraints to avoid post-hoc metric selection bias.

\textbf{Depth guidance (advanced).} At the advanced level, the same concept should be treated as a research variable: students should formulate competing hypotheses, design controlled ablations, and reason about how uncertainty in one subsystem propagates into frontier movement.

\textbf{Operational checklist.} Before accepting any run result for this module, verify: (i) metric definitions are unchanged across runs, (ii) scenario and seed lists are fixed, (iii) checkpoint metadata is parseable by automated evaluators, (iv) run logs include safety channels, and (v) threshold tier classification is reported together with raw reward and cost.

\begin{itemize}[leftmargin=*]
  \item Principle statement: intervention-invariant features improve robustness.
  \item Primary challenge: spurious visual correlations leak into value estimates.
  \item Mitigation path: introduce data corruption stress tests and track degradation under lighting/camera shifts.
  \item Verification channels: reward, cost, lambda, risk, block rate.
  \item Success criterion: positive frontier trend with stable safety traces.
  \item Documentation rule: report both mean and seed-level outcomes.
\end{itemize}


\subsubsection{Module 9: Hierarchy and Temporal Credit (Foundational View)}
\textbf{Principle.} long-horizon tasks need abstraction to shorten credit chains. In AMASA-Complex v2 this appears as a constraint on both algorithm design and experiment process. The key point is that no isolated formula is sufficient; correctness emerges from interaction between objective, data regime, guard behavior, and reproducibility checks.

\textbf{Failure mode.} A common practical breakdown is meta-policy stale transitions break as low-level changes. This usually manifests as deceptive short-run gains followed by cross-seed instability. In safety-critical tasks, this instability must be interpreted as a primary failure signal even when average reward appears competitive.

\textbf{Mitigation.} The project mitigation is to off-policy relabeling and consistency checks keep hierarchy synchronized. This mitigation is encoded in configuration and code so that it can be repeated consistently across seeds and scenarios. It is also paired with evaluation constraints to avoid post-hoc metric selection bias.

\textbf{Depth guidance (foundational).} At the foundational level, students should be able to restate this concept formally, map it to concrete code paths, and identify the minimum telemetry needed to validate implementation correctness.

\textbf{Operational checklist.} Before accepting any run result for this module, verify: (i) metric definitions are unchanged across runs, (ii) scenario and seed lists are fixed, (iii) checkpoint metadata is parseable by automated evaluators, (iv) run logs include safety channels, and (v) threshold tier classification is reported together with raw reward and cost.

\begin{itemize}[leftmargin=*]
  \item Principle statement: long-horizon tasks need abstraction to shorten credit chains.
  \item Primary challenge: meta-policy stale transitions break as low-level changes.
  \item Mitigation path: off-policy relabeling and consistency checks keep hierarchy synchronized.
  \item Verification channels: reward, cost, lambda, risk, block rate.
  \item Success criterion: positive frontier trend with stable safety traces.
  \item Documentation rule: report both mean and seed-level outcomes.
\end{itemize}


\subsubsection{Module 10: Hierarchy and Temporal Credit (Advanced Research View)}
\textbf{Principle.} long-horizon tasks need abstraction to shorten credit chains. In AMASA-Complex v2 this appears as a constraint on both algorithm design and experiment process. The key point is that no isolated formula is sufficient; correctness emerges from interaction between objective, data regime, guard behavior, and reproducibility checks.

\textbf{Failure mode.} A common practical breakdown is meta-policy stale transitions break as low-level changes. This usually manifests as deceptive short-run gains followed by cross-seed instability. In safety-critical tasks, this instability must be interpreted as a primary failure signal even when average reward appears competitive.

\textbf{Mitigation.} The project mitigation is to off-policy relabeling and consistency checks keep hierarchy synchronized. This mitigation is encoded in configuration and code so that it can be repeated consistently across seeds and scenarios. It is also paired with evaluation constraints to avoid post-hoc metric selection bias.

\textbf{Depth guidance (advanced).} At the advanced level, the same concept should be treated as a research variable: students should formulate competing hypotheses, design controlled ablations, and reason about how uncertainty in one subsystem propagates into frontier movement.

\textbf{Operational checklist.} Before accepting any run result for this module, verify: (i) metric definitions are unchanged across runs, (ii) scenario and seed lists are fixed, (iii) checkpoint metadata is parseable by automated evaluators, (iv) run logs include safety channels, and (v) threshold tier classification is reported together with raw reward and cost.

\begin{itemize}[leftmargin=*]
  \item Principle statement: long-horizon tasks need abstraction to shorten credit chains.
  \item Primary challenge: meta-policy stale transitions break as low-level changes.
  \item Mitigation path: off-policy relabeling and consistency checks keep hierarchy synchronized.
  \item Verification channels: reward, cost, lambda, risk, block rate.
  \item Success criterion: positive frontier trend with stable safety traces.
  \item Documentation rule: report both mean and seed-level outcomes.
\end{itemize}


\subsubsection{Module 11: Safe Offline-to-Online Transfer (Foundational View)}
\textbf{Principle.} warm starting policy weights alone is insufficient. In AMASA-Complex v2 this appears as a constraint on both algorithm design and experiment process. The key point is that no isolated formula is sufficient; correctness emerges from interaction between objective, data regime, guard behavior, and reproducibility checks.

\textbf{Failure mode.} A common practical breakdown is cost critics are poorly calibrated after offline pretraining. This usually manifests as deceptive short-run gains followed by cross-seed instability. In safety-critical tasks, this instability must be interpreted as a primary failure signal even when average reward appears competitive.

\textbf{Mitigation.} The project mitigation is to perform value pre-alignment and staged online adaptation. This mitigation is encoded in configuration and code so that it can be repeated consistently across seeds and scenarios. It is also paired with evaluation constraints to avoid post-hoc metric selection bias.

\textbf{Depth guidance (foundational).} At the foundational level, students should be able to restate this concept formally, map it to concrete code paths, and identify the minimum telemetry needed to validate implementation correctness.

\textbf{Operational checklist.} Before accepting any run result for this module, verify: (i) metric definitions are unchanged across runs, (ii) scenario and seed lists are fixed, (iii) checkpoint metadata is parseable by automated evaluators, (iv) run logs include safety channels, and (v) threshold tier classification is reported together with raw reward and cost.

\begin{itemize}[leftmargin=*]
  \item Principle statement: warm starting policy weights alone is insufficient.
  \item Primary challenge: cost critics are poorly calibrated after offline pretraining.
  \item Mitigation path: perform value pre-alignment and staged online adaptation.
  \item Verification channels: reward, cost, lambda, risk, block rate.
  \item Success criterion: positive frontier trend with stable safety traces.
  \item Documentation rule: report both mean and seed-level outcomes.
\end{itemize}


\subsubsection{Module 12: Safe Offline-to-Online Transfer (Advanced Research View)}
\textbf{Principle.} warm starting policy weights alone is insufficient. In AMASA-Complex v2 this appears as a constraint on both algorithm design and experiment process. The key point is that no isolated formula is sufficient; correctness emerges from interaction between objective, data regime, guard behavior, and reproducibility checks.

\textbf{Failure mode.} A common practical breakdown is cost critics are poorly calibrated after offline pretraining. This usually manifests as deceptive short-run gains followed by cross-seed instability. In safety-critical tasks, this instability must be interpreted as a primary failure signal even when average reward appears competitive.

\textbf{Mitigation.} The project mitigation is to perform value pre-alignment and staged online adaptation. This mitigation is encoded in configuration and code so that it can be repeated consistently across seeds and scenarios. It is also paired with evaluation constraints to avoid post-hoc metric selection bias.

\textbf{Depth guidance (advanced).} At the advanced level, the same concept should be treated as a research variable: students should formulate competing hypotheses, design controlled ablations, and reason about how uncertainty in one subsystem propagates into frontier movement.

\textbf{Operational checklist.} Before accepting any run result for this module, verify: (i) metric definitions are unchanged across runs, (ii) scenario and seed lists are fixed, (iii) checkpoint metadata is parseable by automated evaluators, (iv) run logs include safety channels, and (v) threshold tier classification is reported together with raw reward and cost.

\begin{itemize}[leftmargin=*]
  \item Principle statement: warm starting policy weights alone is insufficient.
  \item Primary challenge: cost critics are poorly calibrated after offline pretraining.
  \item Mitigation path: perform value pre-alignment and staged online adaptation.
  \item Verification channels: reward, cost, lambda, risk, block rate.
  \item Success criterion: positive frontier trend with stable safety traces.
  \item Documentation rule: report both mean and seed-level outcomes.
\end{itemize}


\subsubsection{Module 13: SAC-Lagrangian Dynamics (Foundational View)}
\textbf{Principle.} off-policy constrained updates are sample-efficient. In AMASA-Complex v2 this appears as a constraint on both algorithm design and experiment process. The key point is that no isolated formula is sufficient; correctness emerges from interaction between objective, data regime, guard behavior, and reproducibility checks.

\textbf{Failure mode.} A common practical breakdown is reward critics may dominate cost critics in early replay. This usually manifests as deceptive short-run gains followed by cross-seed instability. In safety-critical tasks, this instability must be interpreted as a primary failure signal even when average reward appears competitive.

\textbf{Mitigation.} The project mitigation is to balance replay mixture and monitor critic-scale mismatch. This mitigation is encoded in configuration and code so that it can be repeated consistently across seeds and scenarios. It is also paired with evaluation constraints to avoid post-hoc metric selection bias.

\textbf{Depth guidance (foundational).} At the foundational level, students should be able to restate this concept formally, map it to concrete code paths, and identify the minimum telemetry needed to validate implementation correctness.

\textbf{Operational checklist.} Before accepting any run result for this module, verify: (i) metric definitions are unchanged across runs, (ii) scenario and seed lists are fixed, (iii) checkpoint metadata is parseable by automated evaluators, (iv) run logs include safety channels, and (v) threshold tier classification is reported together with raw reward and cost.

\begin{itemize}[leftmargin=*]
  \item Principle statement: off-policy constrained updates are sample-efficient.
  \item Primary challenge: reward critics may dominate cost critics in early replay.
  \item Mitigation path: balance replay mixture and monitor critic-scale mismatch.
  \item Verification channels: reward, cost, lambda, risk, block rate.
  \item Success criterion: positive frontier trend with stable safety traces.
  \item Documentation rule: report both mean and seed-level outcomes.
\end{itemize}


\subsubsection{Module 14: SAC-Lagrangian Dynamics (Advanced Research View)}
\textbf{Principle.} off-policy constrained updates are sample-efficient. In AMASA-Complex v2 this appears as a constraint on both algorithm design and experiment process. The key point is that no isolated formula is sufficient; correctness emerges from interaction between objective, data regime, guard behavior, and reproducibility checks.

\textbf{Failure mode.} A common practical breakdown is reward critics may dominate cost critics in early replay. This usually manifests as deceptive short-run gains followed by cross-seed instability. In safety-critical tasks, this instability must be interpreted as a primary failure signal even when average reward appears competitive.

\textbf{Mitigation.} The project mitigation is to balance replay mixture and monitor critic-scale mismatch. This mitigation is encoded in configuration and code so that it can be repeated consistently across seeds and scenarios. It is also paired with evaluation constraints to avoid post-hoc metric selection bias.

\textbf{Depth guidance (advanced).} At the advanced level, the same concept should be treated as a research variable: students should formulate competing hypotheses, design controlled ablations, and reason about how uncertainty in one subsystem propagates into frontier movement.

\textbf{Operational checklist.} Before accepting any run result for this module, verify: (i) metric definitions are unchanged across runs, (ii) scenario and seed lists are fixed, (iii) checkpoint metadata is parseable by automated evaluators, (iv) run logs include safety channels, and (v) threshold tier classification is reported together with raw reward and cost.

\begin{itemize}[leftmargin=*]
  \item Principle statement: off-policy constrained updates are sample-efficient.
  \item Primary challenge: reward critics may dominate cost critics in early replay.
  \item Mitigation path: balance replay mixture and monitor critic-scale mismatch.
  \item Verification channels: reward, cost, lambda, risk, block rate.
  \item Success criterion: positive frontier trend with stable safety traces.
  \item Documentation rule: report both mean and seed-level outcomes.
\end{itemize}


\subsubsection{Module 15: PPO-Lagrangian Dynamics (Foundational View)}
\textbf{Principle.} on-policy updates can be robust to stale replay. In AMASA-Complex v2 this appears as a constraint on both algorithm design and experiment process. The key point is that no isolated formula is sufficient; correctness emerges from interaction between objective, data regime, guard behavior, and reproducibility checks.

\textbf{Failure mode.} A common practical breakdown is high variance in cost advantage causes unstable clipping regimes. This usually manifests as deceptive short-run gains followed by cross-seed instability. In safety-critical tasks, this instability must be interpreted as a primary failure signal even when average reward appears competitive.

\textbf{Mitigation.} The project mitigation is to normalize reward and cost channels separately before aggregation. This mitigation is encoded in configuration and code so that it can be repeated consistently across seeds and scenarios. It is also paired with evaluation constraints to avoid post-hoc metric selection bias.

\textbf{Depth guidance (foundational).} At the foundational level, students should be able to restate this concept formally, map it to concrete code paths, and identify the minimum telemetry needed to validate implementation correctness.

\textbf{Operational checklist.} Before accepting any run result for this module, verify: (i) metric definitions are unchanged across runs, (ii) scenario and seed lists are fixed, (iii) checkpoint metadata is parseable by automated evaluators, (iv) run logs include safety channels, and (v) threshold tier classification is reported together with raw reward and cost.

\begin{itemize}[leftmargin=*]
  \item Principle statement: on-policy updates can be robust to stale replay.
  \item Primary challenge: high variance in cost advantage causes unstable clipping regimes.
  \item Mitigation path: normalize reward and cost channels separately before aggregation.
  \item Verification channels: reward, cost, lambda, risk, block rate.
  \item Success criterion: positive frontier trend with stable safety traces.
  \item Documentation rule: report both mean and seed-level outcomes.
\end{itemize}


\subsubsection{Module 16: PPO-Lagrangian Dynamics (Advanced Research View)}
\textbf{Principle.} on-policy updates can be robust to stale replay. In AMASA-Complex v2 this appears as a constraint on both algorithm design and experiment process. The key point is that no isolated formula is sufficient; correctness emerges from interaction between objective, data regime, guard behavior, and reproducibility checks.

\textbf{Failure mode.} A common practical breakdown is high variance in cost advantage causes unstable clipping regimes. This usually manifests as deceptive short-run gains followed by cross-seed instability. In safety-critical tasks, this instability must be interpreted as a primary failure signal even when average reward appears competitive.

\textbf{Mitigation.} The project mitigation is to normalize reward and cost channels separately before aggregation. This mitigation is encoded in configuration and code so that it can be repeated consistently across seeds and scenarios. It is also paired with evaluation constraints to avoid post-hoc metric selection bias.

\textbf{Depth guidance (advanced).} At the advanced level, the same concept should be treated as a research variable: students should formulate competing hypotheses, design controlled ablations, and reason about how uncertainty in one subsystem propagates into frontier movement.

\textbf{Operational checklist.} Before accepting any run result for this module, verify: (i) metric definitions are unchanged across runs, (ii) scenario and seed lists are fixed, (iii) checkpoint metadata is parseable by automated evaluators, (iv) run logs include safety channels, and (v) threshold tier classification is reported together with raw reward and cost.

\begin{itemize}[leftmargin=*]
  \item Principle statement: on-policy updates can be robust to stale replay.
  \item Primary challenge: high variance in cost advantage causes unstable clipping regimes.
  \item Mitigation path: normalize reward and cost channels separately before aggregation.
  \item Verification channels: reward, cost, lambda, risk, block rate.
  \item Success criterion: positive frontier trend with stable safety traces.
  \item Documentation rule: report both mean and seed-level outcomes.
\end{itemize}


\subsubsection{Module 17: Risk Critic Calibration (Foundational View)}
\textbf{Principle.} binary risk scoring supports fast runtime gating. In AMASA-Complex v2 this appears as a constraint on both algorithm design and experiment process. The key point is that no isolated formula is sufficient; correctness emerges from interaction between objective, data regime, guard behavior, and reproducibility checks.

\textbf{Failure mode.} A common practical breakdown is label noise from near-threshold events blurs decision boundaries. This usually manifests as deceptive short-run gains followed by cross-seed instability. In safety-critical tasks, this instability must be interpreted as a primary failure signal even when average reward appears competitive.

\textbf{Mitigation.} The project mitigation is to use short history vectors and periodic calibration checks. This mitigation is encoded in configuration and code so that it can be repeated consistently across seeds and scenarios. It is also paired with evaluation constraints to avoid post-hoc metric selection bias.

\textbf{Depth guidance (foundational).} At the foundational level, students should be able to restate this concept formally, map it to concrete code paths, and identify the minimum telemetry needed to validate implementation correctness.

\textbf{Operational checklist.} Before accepting any run result for this module, verify: (i) metric definitions are unchanged across runs, (ii) scenario and seed lists are fixed, (iii) checkpoint metadata is parseable by automated evaluators, (iv) run logs include safety channels, and (v) threshold tier classification is reported together with raw reward and cost.

\begin{itemize}[leftmargin=*]
  \item Principle statement: binary risk scoring supports fast runtime gating.
  \item Primary challenge: label noise from near-threshold events blurs decision boundaries.
  \item Mitigation path: use short history vectors and periodic calibration checks.
  \item Verification channels: reward, cost, lambda, risk, block rate.
  \item Success criterion: positive frontier trend with stable safety traces.
  \item Documentation rule: report both mean and seed-level outcomes.
\end{itemize}


\subsubsection{Module 18: Risk Critic Calibration (Advanced Research View)}
\textbf{Principle.} binary risk scoring supports fast runtime gating. In AMASA-Complex v2 this appears as a constraint on both algorithm design and experiment process. The key point is that no isolated formula is sufficient; correctness emerges from interaction between objective, data regime, guard behavior, and reproducibility checks.

\textbf{Failure mode.} A common practical breakdown is label noise from near-threshold events blurs decision boundaries. This usually manifests as deceptive short-run gains followed by cross-seed instability. In safety-critical tasks, this instability must be interpreted as a primary failure signal even when average reward appears competitive.

\textbf{Mitigation.} The project mitigation is to use short history vectors and periodic calibration checks. This mitigation is encoded in configuration and code so that it can be repeated consistently across seeds and scenarios. It is also paired with evaluation constraints to avoid post-hoc metric selection bias.

\textbf{Depth guidance (advanced).} At the advanced level, the same concept should be treated as a research variable: students should formulate competing hypotheses, design controlled ablations, and reason about how uncertainty in one subsystem propagates into frontier movement.

\textbf{Operational checklist.} Before accepting any run result for this module, verify: (i) metric definitions are unchanged across runs, (ii) scenario and seed lists are fixed, (iii) checkpoint metadata is parseable by automated evaluators, (iv) run logs include safety channels, and (v) threshold tier classification is reported together with raw reward and cost.

\begin{itemize}[leftmargin=*]
  \item Principle statement: binary risk scoring supports fast runtime gating.
  \item Primary challenge: label noise from near-threshold events blurs decision boundaries.
  \item Mitigation path: use short history vectors and periodic calibration checks.
  \item Verification channels: reward, cost, lambda, risk, block rate.
  \item Success criterion: positive frontier trend with stable safety traces.
  \item Documentation rule: report both mean and seed-level outcomes.
\end{itemize}


\subsubsection{Module 19: Decision Tree Shield Interpretability (Foundational View)}
\textbf{Principle.} rule-based filtering provides human-readable explanations. In AMASA-Complex v2 this appears as a constraint on both algorithm design and experiment process. The key point is that no isolated formula is sufficient; correctness emerges from interaction between objective, data regime, guard behavior, and reproducibility checks.

\textbf{Failure mode.} A common practical breakdown is overly strict rules create conservative deadlocks. This usually manifests as deceptive short-run gains followed by cross-seed instability. In safety-critical tasks, this instability must be interpreted as a primary failure signal even when average reward appears competitive.

\textbf{Mitigation.} The project mitigation is to periodically retrain shield from updated safe/unsafe labels. This mitigation is encoded in configuration and code so that it can be repeated consistently across seeds and scenarios. It is also paired with evaluation constraints to avoid post-hoc metric selection bias.

\textbf{Depth guidance (foundational).} At the foundational level, students should be able to restate this concept formally, map it to concrete code paths, and identify the minimum telemetry needed to validate implementation correctness.

\textbf{Operational checklist.} Before accepting any run result for this module, verify: (i) metric definitions are unchanged across runs, (ii) scenario and seed lists are fixed, (iii) checkpoint metadata is parseable by automated evaluators, (iv) run logs include safety channels, and (v) threshold tier classification is reported together with raw reward and cost.

\begin{itemize}[leftmargin=*]
  \item Principle statement: rule-based filtering provides human-readable explanations.
  \item Primary challenge: overly strict rules create conservative deadlocks.
  \item Mitigation path: periodically retrain shield from updated safe/unsafe labels.
  \item Verification channels: reward, cost, lambda, risk, block rate.
  \item Success criterion: positive frontier trend with stable safety traces.
  \item Documentation rule: report both mean and seed-level outcomes.
\end{itemize}


\subsubsection{Module 20: Decision Tree Shield Interpretability (Advanced Research View)}
\textbf{Principle.} rule-based filtering provides human-readable explanations. In AMASA-Complex v2 this appears as a constraint on both algorithm design and experiment process. The key point is that no isolated formula is sufficient; correctness emerges from interaction between objective, data regime, guard behavior, and reproducibility checks.

\textbf{Failure mode.} A common practical breakdown is overly strict rules create conservative deadlocks. This usually manifests as deceptive short-run gains followed by cross-seed instability. In safety-critical tasks, this instability must be interpreted as a primary failure signal even when average reward appears competitive.

\textbf{Mitigation.} The project mitigation is to periodically retrain shield from updated safe/unsafe labels. This mitigation is encoded in configuration and code so that it can be repeated consistently across seeds and scenarios. It is also paired with evaluation constraints to avoid post-hoc metric selection bias.

\textbf{Depth guidance (advanced).} At the advanced level, the same concept should be treated as a research variable: students should formulate competing hypotheses, design controlled ablations, and reason about how uncertainty in one subsystem propagates into frontier movement.

\textbf{Operational checklist.} Before accepting any run result for this module, verify: (i) metric definitions are unchanged across runs, (ii) scenario and seed lists are fixed, (iii) checkpoint metadata is parseable by automated evaluators, (iv) run logs include safety channels, and (v) threshold tier classification is reported together with raw reward and cost.

\begin{itemize}[leftmargin=*]
  \item Principle statement: rule-based filtering provides human-readable explanations.
  \item Primary challenge: overly strict rules create conservative deadlocks.
  \item Mitigation path: periodically retrain shield from updated safe/unsafe labels.
  \item Verification channels: reward, cost, lambda, risk, block rate.
  \item Success criterion: positive frontier trend with stable safety traces.
  \item Documentation rule: report both mean and seed-level outcomes.
\end{itemize}


\subsubsection{Module 21: PID Gain Scheduling (Foundational View)}
\textbf{Principle.} gain choice controls responsiveness and overshoot. In AMASA-Complex v2 this appears as a constraint on both algorithm design and experiment process. The key point is that no isolated formula is sufficient; correctness emerges from interaction between objective, data regime, guard behavior, and reproducibility checks.

\textbf{Failure mode.} A common practical breakdown is high kp can lock policy in high-lambda mode. This usually manifests as deceptive short-run gains followed by cross-seed instability. In safety-critical tasks, this instability must be interpreted as a primary failure signal even when average reward appears competitive.

\textbf{Mitigation.} The project mitigation is to sweep kp and kd jointly and inspect frontier score. This mitigation is encoded in configuration and code so that it can be repeated consistently across seeds and scenarios. It is also paired with evaluation constraints to avoid post-hoc metric selection bias.

\textbf{Depth guidance (foundational).} At the foundational level, students should be able to restate this concept formally, map it to concrete code paths, and identify the minimum telemetry needed to validate implementation correctness.

\textbf{Operational checklist.} Before accepting any run result for this module, verify: (i) metric definitions are unchanged across runs, (ii) scenario and seed lists are fixed, (iii) checkpoint metadata is parseable by automated evaluators, (iv) run logs include safety channels, and (v) threshold tier classification is reported together with raw reward and cost.

\begin{itemize}[leftmargin=*]
  \item Principle statement: gain choice controls responsiveness and overshoot.
  \item Primary challenge: high kp can lock policy in high-lambda mode.
  \item Mitigation path: sweep kp and kd jointly and inspect frontier score.
  \item Verification channels: reward, cost, lambda, risk, block rate.
  \item Success criterion: positive frontier trend with stable safety traces.
  \item Documentation rule: report both mean and seed-level outcomes.
\end{itemize}


\subsubsection{Module 22: PID Gain Scheduling (Advanced Research View)}
\textbf{Principle.} gain choice controls responsiveness and overshoot. In AMASA-Complex v2 this appears as a constraint on both algorithm design and experiment process. The key point is that no isolated formula is sufficient; correctness emerges from interaction between objective, data regime, guard behavior, and reproducibility checks.

\textbf{Failure mode.} A common practical breakdown is high kp can lock policy in high-lambda mode. This usually manifests as deceptive short-run gains followed by cross-seed instability. In safety-critical tasks, this instability must be interpreted as a primary failure signal even when average reward appears competitive.

\textbf{Mitigation.} The project mitigation is to sweep kp and kd jointly and inspect frontier score. This mitigation is encoded in configuration and code so that it can be repeated consistently across seeds and scenarios. It is also paired with evaluation constraints to avoid post-hoc metric selection bias.

\textbf{Depth guidance (advanced).} At the advanced level, the same concept should be treated as a research variable: students should formulate competing hypotheses, design controlled ablations, and reason about how uncertainty in one subsystem propagates into frontier movement.

\textbf{Operational checklist.} Before accepting any run result for this module, verify: (i) metric definitions are unchanged across runs, (ii) scenario and seed lists are fixed, (iii) checkpoint metadata is parseable by automated evaluators, (iv) run logs include safety channels, and (v) threshold tier classification is reported together with raw reward and cost.

\begin{itemize}[leftmargin=*]
  \item Principle statement: gain choice controls responsiveness and overshoot.
  \item Primary challenge: high kp can lock policy in high-lambda mode.
  \item Mitigation path: sweep kp and kd jointly and inspect frontier score.
  \item Verification channels: reward, cost, lambda, risk, block rate.
  \item Success criterion: positive frontier trend with stable safety traces.
  \item Documentation rule: report both mean and seed-level outcomes.
\end{itemize}


\subsubsection{Module 23: Reward-Cost Pareto Evaluation (Foundational View)}
\textbf{Principle.} single-metric scoring hides safety violations. In AMASA-Complex v2 this appears as a constraint on both algorithm design and experiment process. The key point is that no isolated formula is sufficient; correctness emerges from interaction between objective, data regime, guard behavior, and reproducibility checks.

\textbf{Failure mode.} A common practical breakdown is benchmark ranking changes under different cost penalties. This usually manifests as deceptive short-run gains followed by cross-seed instability. In safety-critical tasks, this instability must be interpreted as a primary failure signal even when average reward appears competitive.

\textbf{Mitigation.} The project mitigation is to report both raw Pareto and weighted frontier scores. This mitigation is encoded in configuration and code so that it can be repeated consistently across seeds and scenarios. It is also paired with evaluation constraints to avoid post-hoc metric selection bias.

\textbf{Depth guidance (foundational).} At the foundational level, students should be able to restate this concept formally, map it to concrete code paths, and identify the minimum telemetry needed to validate implementation correctness.

\textbf{Operational checklist.} Before accepting any run result for this module, verify: (i) metric definitions are unchanged across runs, (ii) scenario and seed lists are fixed, (iii) checkpoint metadata is parseable by automated evaluators, (iv) run logs include safety channels, and (v) threshold tier classification is reported together with raw reward and cost.

\begin{itemize}[leftmargin=*]
  \item Principle statement: single-metric scoring hides safety violations.
  \item Primary challenge: benchmark ranking changes under different cost penalties.
  \item Mitigation path: report both raw Pareto and weighted frontier scores.
  \item Verification channels: reward, cost, lambda, risk, block rate.
  \item Success criterion: positive frontier trend with stable safety traces.
  \item Documentation rule: report both mean and seed-level outcomes.
\end{itemize}


\subsubsection{Module 24: Reward-Cost Pareto Evaluation (Advanced Research View)}
\textbf{Principle.} single-metric scoring hides safety violations. In AMASA-Complex v2 this appears as a constraint on both algorithm design and experiment process. The key point is that no isolated formula is sufficient; correctness emerges from interaction between objective, data regime, guard behavior, and reproducibility checks.

\textbf{Failure mode.} A common practical breakdown is benchmark ranking changes under different cost penalties. This usually manifests as deceptive short-run gains followed by cross-seed instability. In safety-critical tasks, this instability must be interpreted as a primary failure signal even when average reward appears competitive.

\textbf{Mitigation.} The project mitigation is to report both raw Pareto and weighted frontier scores. This mitigation is encoded in configuration and code so that it can be repeated consistently across seeds and scenarios. It is also paired with evaluation constraints to avoid post-hoc metric selection bias.

\textbf{Depth guidance (advanced).} At the advanced level, the same concept should be treated as a research variable: students should formulate competing hypotheses, design controlled ablations, and reason about how uncertainty in one subsystem propagates into frontier movement.

\textbf{Operational checklist.} Before accepting any run result for this module, verify: (i) metric definitions are unchanged across runs, (ii) scenario and seed lists are fixed, (iii) checkpoint metadata is parseable by automated evaluators, (iv) run logs include safety channels, and (v) threshold tier classification is reported together with raw reward and cost.

\begin{itemize}[leftmargin=*]
  \item Principle statement: single-metric scoring hides safety violations.
  \item Primary challenge: benchmark ranking changes under different cost penalties.
  \item Mitigation path: report both raw Pareto and weighted frontier scores.
  \item Verification channels: reward, cost, lambda, risk, block rate.
  \item Success criterion: positive frontier trend with stable safety traces.
  \item Documentation rule: report both mean and seed-level outcomes.
\end{itemize}


\subsubsection{Module 25: Scenario Randomization (Foundational View)}
\textbf{Principle.} domain randomization improves transfer realism. In AMASA-Complex v2 this appears as a constraint on both algorithm design and experiment process. The key point is that no isolated formula is sufficient; correctness emerges from interaction between objective, data regime, guard behavior, and reproducibility checks.

\textbf{Failure mode.} A common practical breakdown is aggressive corruption can collapse exploration. This usually manifests as deceptive short-run gains followed by cross-seed instability. In safety-critical tasks, this instability must be interpreted as a primary failure signal even when average reward appears competitive.

\textbf{Mitigation.} The project mitigation is to mix nominal and perturbed episodes through curriculum. This mitigation is encoded in configuration and code so that it can be repeated consistently across seeds and scenarios. It is also paired with evaluation constraints to avoid post-hoc metric selection bias.

\textbf{Depth guidance (foundational).} At the foundational level, students should be able to restate this concept formally, map it to concrete code paths, and identify the minimum telemetry needed to validate implementation correctness.

\textbf{Operational checklist.} Before accepting any run result for this module, verify: (i) metric definitions are unchanged across runs, (ii) scenario and seed lists are fixed, (iii) checkpoint metadata is parseable by automated evaluators, (iv) run logs include safety channels, and (v) threshold tier classification is reported together with raw reward and cost.

\begin{itemize}[leftmargin=*]
  \item Principle statement: domain randomization improves transfer realism.
  \item Primary challenge: aggressive corruption can collapse exploration.
  \item Mitigation path: mix nominal and perturbed episodes through curriculum.
  \item Verification channels: reward, cost, lambda, risk, block rate.
  \item Success criterion: positive frontier trend with stable safety traces.
  \item Documentation rule: report both mean and seed-level outcomes.
\end{itemize}


\subsubsection{Module 26: Scenario Randomization (Advanced Research View)}
\textbf{Principle.} domain randomization improves transfer realism. In AMASA-Complex v2 this appears as a constraint on both algorithm design and experiment process. The key point is that no isolated formula is sufficient; correctness emerges from interaction between objective, data regime, guard behavior, and reproducibility checks.

\textbf{Failure mode.} A common practical breakdown is aggressive corruption can collapse exploration. This usually manifests as deceptive short-run gains followed by cross-seed instability. In safety-critical tasks, this instability must be interpreted as a primary failure signal even when average reward appears competitive.

\textbf{Mitigation.} The project mitigation is to mix nominal and perturbed episodes through curriculum. This mitigation is encoded in configuration and code so that it can be repeated consistently across seeds and scenarios. It is also paired with evaluation constraints to avoid post-hoc metric selection bias.

\textbf{Depth guidance (advanced).} At the advanced level, the same concept should be treated as a research variable: students should formulate competing hypotheses, design controlled ablations, and reason about how uncertainty in one subsystem propagates into frontier movement.

\textbf{Operational checklist.} Before accepting any run result for this module, verify: (i) metric definitions are unchanged across runs, (ii) scenario and seed lists are fixed, (iii) checkpoint metadata is parseable by automated evaluators, (iv) run logs include safety channels, and (v) threshold tier classification is reported together with raw reward and cost.

\begin{itemize}[leftmargin=*]
  \item Principle statement: domain randomization improves transfer realism.
  \item Primary challenge: aggressive corruption can collapse exploration.
  \item Mitigation path: mix nominal and perturbed episodes through curriculum.
  \item Verification channels: reward, cost, lambda, risk, block rate.
  \item Success criterion: positive frontier trend with stable safety traces.
  \item Documentation rule: report both mean and seed-level outcomes.
\end{itemize}


\subsubsection{Module 27: Evaluation Reproducibility (Foundational View)}
\textbf{Principle.} config-first experiments reduce accidental drift. In AMASA-Complex v2 this appears as a constraint on both algorithm design and experiment process. The key point is that no isolated formula is sufficient; correctness emerges from interaction between objective, data regime, guard behavior, and reproducibility checks.

\textbf{Failure mode.} A common practical breakdown is overlay precedence bugs silently alter run budgets. This usually manifests as deceptive short-run gains followed by cross-seed instability. In safety-critical tasks, this instability must be interpreted as a primary failure signal even when average reward appears competitive.

\textbf{Mitigation.} The project mitigation is to enforce deterministic merge order with schema checks. This mitigation is encoded in configuration and code so that it can be repeated consistently across seeds and scenarios. It is also paired with evaluation constraints to avoid post-hoc metric selection bias.

\textbf{Depth guidance (foundational).} At the foundational level, students should be able to restate this concept formally, map it to concrete code paths, and identify the minimum telemetry needed to validate implementation correctness.

\textbf{Operational checklist.} Before accepting any run result for this module, verify: (i) metric definitions are unchanged across runs, (ii) scenario and seed lists are fixed, (iii) checkpoint metadata is parseable by automated evaluators, (iv) run logs include safety channels, and (v) threshold tier classification is reported together with raw reward and cost.

\begin{itemize}[leftmargin=*]
  \item Principle statement: config-first experiments reduce accidental drift.
  \item Primary challenge: overlay precedence bugs silently alter run budgets.
  \item Mitigation path: enforce deterministic merge order with schema checks.
  \item Verification channels: reward, cost, lambda, risk, block rate.
  \item Success criterion: positive frontier trend with stable safety traces.
  \item Documentation rule: report both mean and seed-level outcomes.
\end{itemize}


\subsubsection{Module 28: Evaluation Reproducibility (Advanced Research View)}
\textbf{Principle.} config-first experiments reduce accidental drift. In AMASA-Complex v2 this appears as a constraint on both algorithm design and experiment process. The key point is that no isolated formula is sufficient; correctness emerges from interaction between objective, data regime, guard behavior, and reproducibility checks.

\textbf{Failure mode.} A common practical breakdown is overlay precedence bugs silently alter run budgets. This usually manifests as deceptive short-run gains followed by cross-seed instability. In safety-critical tasks, this instability must be interpreted as a primary failure signal even when average reward appears competitive.

\textbf{Mitigation.} The project mitigation is to enforce deterministic merge order with schema checks. This mitigation is encoded in configuration and code so that it can be repeated consistently across seeds and scenarios. It is also paired with evaluation constraints to avoid post-hoc metric selection bias.

\textbf{Depth guidance (advanced).} At the advanced level, the same concept should be treated as a research variable: students should formulate competing hypotheses, design controlled ablations, and reason about how uncertainty in one subsystem propagates into frontier movement.

\textbf{Operational checklist.} Before accepting any run result for this module, verify: (i) metric definitions are unchanged across runs, (ii) scenario and seed lists are fixed, (iii) checkpoint metadata is parseable by automated evaluators, (iv) run logs include safety channels, and (v) threshold tier classification is reported together with raw reward and cost.

\begin{itemize}[leftmargin=*]
  \item Principle statement: config-first experiments reduce accidental drift.
  \item Primary challenge: overlay precedence bugs silently alter run budgets.
  \item Mitigation path: enforce deterministic merge order with schema checks.
  \item Verification channels: reward, cost, lambda, risk, block rate.
  \item Success criterion: positive frontier trend with stable safety traces.
  \item Documentation rule: report both mean and seed-level outcomes.
\end{itemize}


\subsubsection{Module 29: Checkpoint Semantics (Foundational View)}
\textbf{Principle.} metadata consistency is required for automated evaluation. In AMASA-Complex v2 this appears as a constraint on both algorithm design and experiment process. The key point is that no isolated formula is sufficient; correctness emerges from interaction between objective, data regime, guard behavior, and reproducibility checks.

\textbf{Failure mode.} A common practical breakdown is heterogeneous checkpoint keys break loaders. This usually manifests as deceptive short-run gains followed by cross-seed instability. In safety-critical tasks, this instability must be interpreted as a primary failure signal even when average reward appears competitive.

\textbf{Mitigation.} The project mitigation is to standardize algo tags and provide fallback serializers. This mitigation is encoded in configuration and code so that it can be repeated consistently across seeds and scenarios. It is also paired with evaluation constraints to avoid post-hoc metric selection bias.

\textbf{Depth guidance (foundational).} At the foundational level, students should be able to restate this concept formally, map it to concrete code paths, and identify the minimum telemetry needed to validate implementation correctness.

\textbf{Operational checklist.} Before accepting any run result for this module, verify: (i) metric definitions are unchanged across runs, (ii) scenario and seed lists are fixed, (iii) checkpoint metadata is parseable by automated evaluators, (iv) run logs include safety channels, and (v) threshold tier classification is reported together with raw reward and cost.

\begin{itemize}[leftmargin=*]
  \item Principle statement: metadata consistency is required for automated evaluation.
  \item Primary challenge: heterogeneous checkpoint keys break loaders.
  \item Mitigation path: standardize algo tags and provide fallback serializers.
  \item Verification channels: reward, cost, lambda, risk, block rate.
  \item Success criterion: positive frontier trend with stable safety traces.
  \item Documentation rule: report both mean and seed-level outcomes.
\end{itemize}


\subsubsection{Module 30: Checkpoint Semantics (Advanced Research View)}
\textbf{Principle.} metadata consistency is required for automated evaluation. In AMASA-Complex v2 this appears as a constraint on both algorithm design and experiment process. The key point is that no isolated formula is sufficient; correctness emerges from interaction between objective, data regime, guard behavior, and reproducibility checks.

\textbf{Failure mode.} A common practical breakdown is heterogeneous checkpoint keys break loaders. This usually manifests as deceptive short-run gains followed by cross-seed instability. In safety-critical tasks, this instability must be interpreted as a primary failure signal even when average reward appears competitive.

\textbf{Mitigation.} The project mitigation is to standardize algo tags and provide fallback serializers. This mitigation is encoded in configuration and code so that it can be repeated consistently across seeds and scenarios. It is also paired with evaluation constraints to avoid post-hoc metric selection bias.

\textbf{Depth guidance (advanced).} At the advanced level, the same concept should be treated as a research variable: students should formulate competing hypotheses, design controlled ablations, and reason about how uncertainty in one subsystem propagates into frontier movement.

\textbf{Operational checklist.} Before accepting any run result for this module, verify: (i) metric definitions are unchanged across runs, (ii) scenario and seed lists are fixed, (iii) checkpoint metadata is parseable by automated evaluators, (iv) run logs include safety channels, and (v) threshold tier classification is reported together with raw reward and cost.

\begin{itemize}[leftmargin=*]
  \item Principle statement: metadata consistency is required for automated evaluation.
  \item Primary challenge: heterogeneous checkpoint keys break loaders.
  \item Mitigation path: standardize algo tags and provide fallback serializers.
  \item Verification channels: reward, cost, lambda, risk, block rate.
  \item Success criterion: positive frontier trend with stable safety traces.
  \item Documentation rule: report both mean and seed-level outcomes.
\end{itemize}


\subsubsection{Module 31: Safety Telemetry Engineering (Foundational View)}
\textbf{Principle.} step-level logs reveal hidden instability. In AMASA-Complex v2 this appears as a constraint on both algorithm design and experiment process. The key point is that no isolated formula is sufficient; correctness emerges from interaction between objective, data regime, guard behavior, and reproducibility checks.

\textbf{Failure mode.} A common practical breakdown is episode-only metrics hide transient safety spikes. This usually manifests as deceptive short-run gains followed by cross-seed instability. In safety-critical tasks, this instability must be interpreted as a primary failure signal even when average reward appears competitive.

\textbf{Mitigation.} The project mitigation is to export lambda, risk, block rate, force, and corridor traces. This mitigation is encoded in configuration and code so that it can be repeated consistently across seeds and scenarios. It is also paired with evaluation constraints to avoid post-hoc metric selection bias.

\textbf{Depth guidance (foundational).} At the foundational level, students should be able to restate this concept formally, map it to concrete code paths, and identify the minimum telemetry needed to validate implementation correctness.

\textbf{Operational checklist.} Before accepting any run result for this module, verify: (i) metric definitions are unchanged across runs, (ii) scenario and seed lists are fixed, (iii) checkpoint metadata is parseable by automated evaluators, (iv) run logs include safety channels, and (v) threshold tier classification is reported together with raw reward and cost.

\begin{itemize}[leftmargin=*]
  \item Principle statement: step-level logs reveal hidden instability.
  \item Primary challenge: episode-only metrics hide transient safety spikes.
  \item Mitigation path: export lambda, risk, block rate, force, and corridor traces.
  \item Verification channels: reward, cost, lambda, risk, block rate.
  \item Success criterion: positive frontier trend with stable safety traces.
  \item Documentation rule: report both mean and seed-level outcomes.
\end{itemize}


\subsubsection{Module 32: Safety Telemetry Engineering (Advanced Research View)}
\textbf{Principle.} step-level logs reveal hidden instability. In AMASA-Complex v2 this appears as a constraint on both algorithm design and experiment process. The key point is that no isolated formula is sufficient; correctness emerges from interaction between objective, data regime, guard behavior, and reproducibility checks.

\textbf{Failure mode.} A common practical breakdown is episode-only metrics hide transient safety spikes. This usually manifests as deceptive short-run gains followed by cross-seed instability. In safety-critical tasks, this instability must be interpreted as a primary failure signal even when average reward appears competitive.

\textbf{Mitigation.} The project mitigation is to export lambda, risk, block rate, force, and corridor traces. This mitigation is encoded in configuration and code so that it can be repeated consistently across seeds and scenarios. It is also paired with evaluation constraints to avoid post-hoc metric selection bias.

\textbf{Depth guidance (advanced).} At the advanced level, the same concept should be treated as a research variable: students should formulate competing hypotheses, design controlled ablations, and reason about how uncertainty in one subsystem propagates into frontier movement.

\textbf{Operational checklist.} Before accepting any run result for this module, verify: (i) metric definitions are unchanged across runs, (ii) scenario and seed lists are fixed, (iii) checkpoint metadata is parseable by automated evaluators, (iv) run logs include safety channels, and (v) threshold tier classification is reported together with raw reward and cost.

\begin{itemize}[leftmargin=*]
  \item Principle statement: step-level logs reveal hidden instability.
  \item Primary challenge: episode-only metrics hide transient safety spikes.
  \item Mitigation path: export lambda, risk, block rate, force, and corridor traces.
  \item Verification channels: reward, cost, lambda, risk, block rate.
  \item Success criterion: positive frontier trend with stable safety traces.
  \item Documentation rule: report both mean and seed-level outcomes.
\end{itemize}


\subsubsection{Module 33: Curriculum Design (Foundational View)}
\textbf{Principle.} phase scheduling improves stability under complexity. In AMASA-Complex v2 this appears as a constraint on both algorithm design and experiment process. The key point is that no isolated formula is sufficient; correctness emerges from interaction between objective, data regime, guard behavior, and reproducibility checks.

\textbf{Failure mode.} A common practical breakdown is jumping directly to adversarial regimes causes collapse. This usually manifests as deceptive short-run gains followed by cross-seed instability. In safety-critical tasks, this instability must be interpreted as a primary failure signal even when average reward appears competitive.

\textbf{Mitigation.} The project mitigation is to progress from nominal to perturbed to adversarial mixtures. This mitigation is encoded in configuration and code so that it can be repeated consistently across seeds and scenarios. It is also paired with evaluation constraints to avoid post-hoc metric selection bias.

\textbf{Depth guidance (foundational).} At the foundational level, students should be able to restate this concept formally, map it to concrete code paths, and identify the minimum telemetry needed to validate implementation correctness.

\textbf{Operational checklist.} Before accepting any run result for this module, verify: (i) metric definitions are unchanged across runs, (ii) scenario and seed lists are fixed, (iii) checkpoint metadata is parseable by automated evaluators, (iv) run logs include safety channels, and (v) threshold tier classification is reported together with raw reward and cost.

\begin{itemize}[leftmargin=*]
  \item Principle statement: phase scheduling improves stability under complexity.
  \item Primary challenge: jumping directly to adversarial regimes causes collapse.
  \item Mitigation path: progress from nominal to perturbed to adversarial mixtures.
  \item Verification channels: reward, cost, lambda, risk, block rate.
  \item Success criterion: positive frontier trend with stable safety traces.
  \item Documentation rule: report both mean and seed-level outcomes.
\end{itemize}


\subsubsection{Module 34: Curriculum Design (Advanced Research View)}
\textbf{Principle.} phase scheduling improves stability under complexity. In AMASA-Complex v2 this appears as a constraint on both algorithm design and experiment process. The key point is that no isolated formula is sufficient; correctness emerges from interaction between objective, data regime, guard behavior, and reproducibility checks.

\textbf{Failure mode.} A common practical breakdown is jumping directly to adversarial regimes causes collapse. This usually manifests as deceptive short-run gains followed by cross-seed instability. In safety-critical tasks, this instability must be interpreted as a primary failure signal even when average reward appears competitive.

\textbf{Mitigation.} The project mitigation is to progress from nominal to perturbed to adversarial mixtures. This mitigation is encoded in configuration and code so that it can be repeated consistently across seeds and scenarios. It is also paired with evaluation constraints to avoid post-hoc metric selection bias.

\textbf{Depth guidance (advanced).} At the advanced level, the same concept should be treated as a research variable: students should formulate competing hypotheses, design controlled ablations, and reason about how uncertainty in one subsystem propagates into frontier movement.

\textbf{Operational checklist.} Before accepting any run result for this module, verify: (i) metric definitions are unchanged across runs, (ii) scenario and seed lists are fixed, (iii) checkpoint metadata is parseable by automated evaluators, (iv) run logs include safety channels, and (v) threshold tier classification is reported together with raw reward and cost.

\begin{itemize}[leftmargin=*]
  \item Principle statement: phase scheduling improves stability under complexity.
  \item Primary challenge: jumping directly to adversarial regimes causes collapse.
  \item Mitigation path: progress from nominal to perturbed to adversarial mixtures.
  \item Verification channels: reward, cost, lambda, risk, block rate.
  \item Success criterion: positive frontier trend with stable safety traces.
  \item Documentation rule: report both mean and seed-level outcomes.
\end{itemize}


\subsubsection{Module 35: Acceptance Thresholds (Foundational View)}
\textbf{Principle.} tiered thresholds clarify minimum/target/stretch outcomes. In AMASA-Complex v2 this appears as a constraint on both algorithm design and experiment process. The key point is that no isolated formula is sufficient; correctness emerges from interaction between objective, data regime, guard behavior, and reproducibility checks.

\textbf{Failure mode.} A common practical breakdown is single pass/fail hides meaningful progress. This usually manifests as deceptive short-run gains followed by cross-seed instability. In safety-critical tasks, this instability must be interpreted as a primary failure signal even when average reward appears competitive.

\textbf{Mitigation.} The project mitigation is to track tier movement over retraining cycles. This mitigation is encoded in configuration and code so that it can be repeated consistently across seeds and scenarios. It is also paired with evaluation constraints to avoid post-hoc metric selection bias.

\textbf{Depth guidance (foundational).} At the foundational level, students should be able to restate this concept formally, map it to concrete code paths, and identify the minimum telemetry needed to validate implementation correctness.

\textbf{Operational checklist.} Before accepting any run result for this module, verify: (i) metric definitions are unchanged across runs, (ii) scenario and seed lists are fixed, (iii) checkpoint metadata is parseable by automated evaluators, (iv) run logs include safety channels, and (v) threshold tier classification is reported together with raw reward and cost.

\begin{itemize}[leftmargin=*]
  \item Principle statement: tiered thresholds clarify minimum/target/stretch outcomes.
  \item Primary challenge: single pass/fail hides meaningful progress.
  \item Mitigation path: track tier movement over retraining cycles.
  \item Verification channels: reward, cost, lambda, risk, block rate.
  \item Success criterion: positive frontier trend with stable safety traces.
  \item Documentation rule: report both mean and seed-level outcomes.
\end{itemize}


\subsubsection{Module 36: Acceptance Thresholds (Advanced Research View)}
\textbf{Principle.} tiered thresholds clarify minimum/target/stretch outcomes. In AMASA-Complex v2 this appears as a constraint on both algorithm design and experiment process. The key point is that no isolated formula is sufficient; correctness emerges from interaction between objective, data regime, guard behavior, and reproducibility checks.

\textbf{Failure mode.} A common practical breakdown is single pass/fail hides meaningful progress. This usually manifests as deceptive short-run gains followed by cross-seed instability. In safety-critical tasks, this instability must be interpreted as a primary failure signal even when average reward appears competitive.

\textbf{Mitigation.} The project mitigation is to track tier movement over retraining cycles. This mitigation is encoded in configuration and code so that it can be repeated consistently across seeds and scenarios. It is also paired with evaluation constraints to avoid post-hoc metric selection bias.

\textbf{Depth guidance (advanced).} At the advanced level, the same concept should be treated as a research variable: students should formulate competing hypotheses, design controlled ablations, and reason about how uncertainty in one subsystem propagates into frontier movement.

\textbf{Operational checklist.} Before accepting any run result for this module, verify: (i) metric definitions are unchanged across runs, (ii) scenario and seed lists are fixed, (iii) checkpoint metadata is parseable by automated evaluators, (iv) run logs include safety channels, and (v) threshold tier classification is reported together with raw reward and cost.

\begin{itemize}[leftmargin=*]
  \item Principle statement: tiered thresholds clarify minimum/target/stretch outcomes.
  \item Primary challenge: single pass/fail hides meaningful progress.
  \item Mitigation path: track tier movement over retraining cycles.
  \item Verification channels: reward, cost, lambda, risk, block rate.
  \item Success criterion: positive frontier trend with stable safety traces.
  \item Documentation rule: report both mean and seed-level outcomes.
\end{itemize}


\subsubsection{Module 37: Failure Taxonomy (Foundational View)}
\textbf{Principle.} structured failure labels accelerate debugging. In AMASA-Complex v2 this appears as a constraint on both algorithm design and experiment process. The key point is that no isolated formula is sufficient; correctness emerges from interaction between objective, data regime, guard behavior, and reproducibility checks.

\textbf{Failure mode.} A common practical breakdown is qualitative logs are hard to compare across seeds. This usually manifests as deceptive short-run gains followed by cross-seed instability. In safety-critical tasks, this instability must be interpreted as a primary failure signal even when average reward appears competitive.

\textbf{Mitigation.} The project mitigation is to define categories: conservative trap, critic drift, shield saturation. This mitigation is encoded in configuration and code so that it can be repeated consistently across seeds and scenarios. It is also paired with evaluation constraints to avoid post-hoc metric selection bias.

\textbf{Depth guidance (foundational).} At the foundational level, students should be able to restate this concept formally, map it to concrete code paths, and identify the minimum telemetry needed to validate implementation correctness.

\textbf{Operational checklist.} Before accepting any run result for this module, verify: (i) metric definitions are unchanged across runs, (ii) scenario and seed lists are fixed, (iii) checkpoint metadata is parseable by automated evaluators, (iv) run logs include safety channels, and (v) threshold tier classification is reported together with raw reward and cost.

\begin{itemize}[leftmargin=*]
  \item Principle statement: structured failure labels accelerate debugging.
  \item Primary challenge: qualitative logs are hard to compare across seeds.
  \item Mitigation path: define categories: conservative trap, critic drift, shield saturation.
  \item Verification channels: reward, cost, lambda, risk, block rate.
  \item Success criterion: positive frontier trend with stable safety traces.
  \item Documentation rule: report both mean and seed-level outcomes.
\end{itemize}


\subsubsection{Module 38: Failure Taxonomy (Advanced Research View)}
\textbf{Principle.} structured failure labels accelerate debugging. In AMASA-Complex v2 this appears as a constraint on both algorithm design and experiment process. The key point is that no isolated formula is sufficient; correctness emerges from interaction between objective, data regime, guard behavior, and reproducibility checks.

\textbf{Failure mode.} A common practical breakdown is qualitative logs are hard to compare across seeds. This usually manifests as deceptive short-run gains followed by cross-seed instability. In safety-critical tasks, this instability must be interpreted as a primary failure signal even when average reward appears competitive.

\textbf{Mitigation.} The project mitigation is to define categories: conservative trap, critic drift, shield saturation. This mitigation is encoded in configuration and code so that it can be repeated consistently across seeds and scenarios. It is also paired with evaluation constraints to avoid post-hoc metric selection bias.

\textbf{Depth guidance (advanced).} At the advanced level, the same concept should be treated as a research variable: students should formulate competing hypotheses, design controlled ablations, and reason about how uncertainty in one subsystem propagates into frontier movement.

\textbf{Operational checklist.} Before accepting any run result for this module, verify: (i) metric definitions are unchanged across runs, (ii) scenario and seed lists are fixed, (iii) checkpoint metadata is parseable by automated evaluators, (iv) run logs include safety channels, and (v) threshold tier classification is reported together with raw reward and cost.

\begin{itemize}[leftmargin=*]
  \item Principle statement: structured failure labels accelerate debugging.
  \item Primary challenge: qualitative logs are hard to compare across seeds.
  \item Mitigation path: define categories: conservative trap, critic drift, shield saturation.
  \item Verification channels: reward, cost, lambda, risk, block rate.
  \item Success criterion: positive frontier trend with stable safety traces.
  \item Documentation rule: report both mean and seed-level outcomes.
\end{itemize}


\subsubsection{Module 39: Ablation Methodology (Foundational View)}
\textbf{Principle.} remove one mechanism at a time for causal attribution. In AMASA-Complex v2 this appears as a constraint on both algorithm design and experiment process. The key point is that no isolated formula is sufficient; correctness emerges from interaction between objective, data regime, guard behavior, and reproducibility checks.

\textbf{Failure mode.} A common practical breakdown is multi-change experiments confound interpretation. This usually manifests as deceptive short-run gains followed by cross-seed instability. In safety-critical tasks, this instability must be interpreted as a primary failure signal even when average reward appears competitive.

\textbf{Mitigation.} The project mitigation is to keep seed set and scenario mix fixed during ablation. This mitigation is encoded in configuration and code so that it can be repeated consistently across seeds and scenarios. It is also paired with evaluation constraints to avoid post-hoc metric selection bias.

\textbf{Depth guidance (foundational).} At the foundational level, students should be able to restate this concept formally, map it to concrete code paths, and identify the minimum telemetry needed to validate implementation correctness.

\textbf{Operational checklist.} Before accepting any run result for this module, verify: (i) metric definitions are unchanged across runs, (ii) scenario and seed lists are fixed, (iii) checkpoint metadata is parseable by automated evaluators, (iv) run logs include safety channels, and (v) threshold tier classification is reported together with raw reward and cost.

\begin{itemize}[leftmargin=*]
  \item Principle statement: remove one mechanism at a time for causal attribution.
  \item Primary challenge: multi-change experiments confound interpretation.
  \item Mitigation path: keep seed set and scenario mix fixed during ablation.
  \item Verification channels: reward, cost, lambda, risk, block rate.
  \item Success criterion: positive frontier trend with stable safety traces.
  \item Documentation rule: report both mean and seed-level outcomes.
\end{itemize}


\subsubsection{Module 40: Ablation Methodology (Advanced Research View)}
\textbf{Principle.} remove one mechanism at a time for causal attribution. In AMASA-Complex v2 this appears as a constraint on both algorithm design and experiment process. The key point is that no isolated formula is sufficient; correctness emerges from interaction between objective, data regime, guard behavior, and reproducibility checks.

\textbf{Failure mode.} A common practical breakdown is multi-change experiments confound interpretation. This usually manifests as deceptive short-run gains followed by cross-seed instability. In safety-critical tasks, this instability must be interpreted as a primary failure signal even when average reward appears competitive.

\textbf{Mitigation.} The project mitigation is to keep seed set and scenario mix fixed during ablation. This mitigation is encoded in configuration and code so that it can be repeated consistently across seeds and scenarios. It is also paired with evaluation constraints to avoid post-hoc metric selection bias.

\textbf{Depth guidance (advanced).} At the advanced level, the same concept should be treated as a research variable: students should formulate competing hypotheses, design controlled ablations, and reason about how uncertainty in one subsystem propagates into frontier movement.

\textbf{Operational checklist.} Before accepting any run result for this module, verify: (i) metric definitions are unchanged across runs, (ii) scenario and seed lists are fixed, (iii) checkpoint metadata is parseable by automated evaluators, (iv) run logs include safety channels, and (v) threshold tier classification is reported together with raw reward and cost.

\begin{itemize}[leftmargin=*]
  \item Principle statement: remove one mechanism at a time for causal attribution.
  \item Primary challenge: multi-change experiments confound interpretation.
  \item Mitigation path: keep seed set and scenario mix fixed during ablation.
  \item Verification channels: reward, cost, lambda, risk, block rate.
  \item Success criterion: positive frontier trend with stable safety traces.
  \item Documentation rule: report both mean and seed-level outcomes.
\end{itemize}


\subsubsection{Module 41: CPU-First Optimization (Foundational View)}
\textbf{Principle.} efficient minibatch pipelines keep the project accessible. In AMASA-Complex v2 this appears as a constraint on both algorithm design and experiment process. The key point is that no isolated formula is sufficient; correctness emerges from interaction between objective, data regime, guard behavior, and reproducibility checks.

\textbf{Failure mode.} A common practical breakdown is slow sweeps discourage rigorous validation. This usually manifests as deceptive short-run gains followed by cross-seed instability. In safety-critical tasks, this instability must be interpreted as a primary failure signal even when average reward appears competitive.

\textbf{Mitigation.} The project mitigation is to reuse dataset buffers and skip completed jobs. This mitigation is encoded in configuration and code so that it can be repeated consistently across seeds and scenarios. It is also paired with evaluation constraints to avoid post-hoc metric selection bias.

\textbf{Depth guidance (foundational).} At the foundational level, students should be able to restate this concept formally, map it to concrete code paths, and identify the minimum telemetry needed to validate implementation correctness.

\textbf{Operational checklist.} Before accepting any run result for this module, verify: (i) metric definitions are unchanged across runs, (ii) scenario and seed lists are fixed, (iii) checkpoint metadata is parseable by automated evaluators, (iv) run logs include safety channels, and (v) threshold tier classification is reported together with raw reward and cost.

\begin{itemize}[leftmargin=*]
  \item Principle statement: efficient minibatch pipelines keep the project accessible.
  \item Primary challenge: slow sweeps discourage rigorous validation.
  \item Mitigation path: reuse dataset buffers and skip completed jobs.
  \item Verification channels: reward, cost, lambda, risk, block rate.
  \item Success criterion: positive frontier trend with stable safety traces.
  \item Documentation rule: report both mean and seed-level outcomes.
\end{itemize}


\subsubsection{Module 42: CPU-First Optimization (Advanced Research View)}
\textbf{Principle.} efficient minibatch pipelines keep the project accessible. In AMASA-Complex v2 this appears as a constraint on both algorithm design and experiment process. The key point is that no isolated formula is sufficient; correctness emerges from interaction between objective, data regime, guard behavior, and reproducibility checks.

\textbf{Failure mode.} A common practical breakdown is slow sweeps discourage rigorous validation. This usually manifests as deceptive short-run gains followed by cross-seed instability. In safety-critical tasks, this instability must be interpreted as a primary failure signal even when average reward appears competitive.

\textbf{Mitigation.} The project mitigation is to reuse dataset buffers and skip completed jobs. This mitigation is encoded in configuration and code so that it can be repeated consistently across seeds and scenarios. It is also paired with evaluation constraints to avoid post-hoc metric selection bias.

\textbf{Depth guidance (advanced).} At the advanced level, the same concept should be treated as a research variable: students should formulate competing hypotheses, design controlled ablations, and reason about how uncertainty in one subsystem propagates into frontier movement.

\textbf{Operational checklist.} Before accepting any run result for this module, verify: (i) metric definitions are unchanged across runs, (ii) scenario and seed lists are fixed, (iii) checkpoint metadata is parseable by automated evaluators, (iv) run logs include safety channels, and (v) threshold tier classification is reported together with raw reward and cost.

\begin{itemize}[leftmargin=*]
  \item Principle statement: efficient minibatch pipelines keep the project accessible.
  \item Primary challenge: slow sweeps discourage rigorous validation.
  \item Mitigation path: reuse dataset buffers and skip completed jobs.
  \item Verification channels: reward, cost, lambda, risk, block rate.
  \item Success criterion: positive frontier trend with stable safety traces.
  \item Documentation rule: report both mean and seed-level outcomes.
\end{itemize}


\subsubsection{Module 43: Backward-Compatible Interfaces (Foundational View)}
\textbf{Principle.} legacy CLIs preserve classroom continuity. In AMASA-Complex v2 this appears as a constraint on both algorithm design and experiment process. The key point is that no isolated formula is sufficient; correctness emerges from interaction between objective, data regime, guard behavior, and reproducibility checks.

\textbf{Failure mode.} A common practical breakdown is breaking flags disrupt earlier assignments. This usually manifests as deceptive short-run gains followed by cross-seed instability. In safety-critical tasks, this instability must be interpreted as a primary failure signal even when average reward appears competitive.

\textbf{Mitigation.} The project mitigation is to add optional config flags while keeping defaults intact. This mitigation is encoded in configuration and code so that it can be repeated consistently across seeds and scenarios. It is also paired with evaluation constraints to avoid post-hoc metric selection bias.

\textbf{Depth guidance (foundational).} At the foundational level, students should be able to restate this concept formally, map it to concrete code paths, and identify the minimum telemetry needed to validate implementation correctness.

\textbf{Operational checklist.} Before accepting any run result for this module, verify: (i) metric definitions are unchanged across runs, (ii) scenario and seed lists are fixed, (iii) checkpoint metadata is parseable by automated evaluators, (iv) run logs include safety channels, and (v) threshold tier classification is reported together with raw reward and cost.

\begin{itemize}[leftmargin=*]
  \item Principle statement: legacy CLIs preserve classroom continuity.
  \item Primary challenge: breaking flags disrupt earlier assignments.
  \item Mitigation path: add optional config flags while keeping defaults intact.
  \item Verification channels: reward, cost, lambda, risk, block rate.
  \item Success criterion: positive frontier trend with stable safety traces.
  \item Documentation rule: report both mean and seed-level outcomes.
\end{itemize}


\subsubsection{Module 44: Backward-Compatible Interfaces (Advanced Research View)}
\textbf{Principle.} legacy CLIs preserve classroom continuity. In AMASA-Complex v2 this appears as a constraint on both algorithm design and experiment process. The key point is that no isolated formula is sufficient; correctness emerges from interaction between objective, data regime, guard behavior, and reproducibility checks.

\textbf{Failure mode.} A common practical breakdown is breaking flags disrupt earlier assignments. This usually manifests as deceptive short-run gains followed by cross-seed instability. In safety-critical tasks, this instability must be interpreted as a primary failure signal even when average reward appears competitive.

\textbf{Mitigation.} The project mitigation is to add optional config flags while keeping defaults intact. This mitigation is encoded in configuration and code so that it can be repeated consistently across seeds and scenarios. It is also paired with evaluation constraints to avoid post-hoc metric selection bias.

\textbf{Depth guidance (advanced).} At the advanced level, the same concept should be treated as a research variable: students should formulate competing hypotheses, design controlled ablations, and reason about how uncertainty in one subsystem propagates into frontier movement.

\textbf{Operational checklist.} Before accepting any run result for this module, verify: (i) metric definitions are unchanged across runs, (ii) scenario and seed lists are fixed, (iii) checkpoint metadata is parseable by automated evaluators, (iv) run logs include safety channels, and (v) threshold tier classification is reported together with raw reward and cost.

\begin{itemize}[leftmargin=*]
  \item Principle statement: legacy CLIs preserve classroom continuity.
  \item Primary challenge: breaking flags disrupt earlier assignments.
  \item Mitigation path: add optional config flags while keeping defaults intact.
  \item Verification channels: reward, cost, lambda, risk, block rate.
  \item Success criterion: positive frontier trend with stable safety traces.
  \item Documentation rule: report both mean and seed-level outcomes.
\end{itemize}


\subsubsection{Module 45: Code Quality and Testing (Foundational View)}
\textbf{Principle.} unit and smoke tests protect refactors. In AMASA-Complex v2 this appears as a constraint on both algorithm design and experiment process. The key point is that no isolated formula is sufficient; correctness emerges from interaction between objective, data regime, guard behavior, and reproducibility checks.

\textbf{Failure mode.} A common practical breakdown is unvalidated merges can silently break training. This usually manifests as deceptive short-run gains followed by cross-seed instability. In safety-critical tasks, this instability must be interpreted as a primary failure signal even when average reward appears competitive.

\textbf{Mitigation.} The project mitigation is to test env determinism, update-step sanity, and CLI smoke paths. This mitigation is encoded in configuration and code so that it can be repeated consistently across seeds and scenarios. It is also paired with evaluation constraints to avoid post-hoc metric selection bias.

\textbf{Depth guidance (foundational).} At the foundational level, students should be able to restate this concept formally, map it to concrete code paths, and identify the minimum telemetry needed to validate implementation correctness.

\textbf{Operational checklist.} Before accepting any run result for this module, verify: (i) metric definitions are unchanged across runs, (ii) scenario and seed lists are fixed, (iii) checkpoint metadata is parseable by automated evaluators, (iv) run logs include safety channels, and (v) threshold tier classification is reported together with raw reward and cost.

\begin{itemize}[leftmargin=*]
  \item Principle statement: unit and smoke tests protect refactors.
  \item Primary challenge: unvalidated merges can silently break training.
  \item Mitigation path: test env determinism, update-step sanity, and CLI smoke paths.
  \item Verification channels: reward, cost, lambda, risk, block rate.
  \item Success criterion: positive frontier trend with stable safety traces.
  \item Documentation rule: report both mean and seed-level outcomes.
\end{itemize}


\subsubsection{Module 46: Code Quality and Testing (Advanced Research View)}
\textbf{Principle.} unit and smoke tests protect refactors. In AMASA-Complex v2 this appears as a constraint on both algorithm design and experiment process. The key point is that no isolated formula is sufficient; correctness emerges from interaction between objective, data regime, guard behavior, and reproducibility checks.

\textbf{Failure mode.} A common practical breakdown is unvalidated merges can silently break training. This usually manifests as deceptive short-run gains followed by cross-seed instability. In safety-critical tasks, this instability must be interpreted as a primary failure signal even when average reward appears competitive.

\textbf{Mitigation.} The project mitigation is to test env determinism, update-step sanity, and CLI smoke paths. This mitigation is encoded in configuration and code so that it can be repeated consistently across seeds and scenarios. It is also paired with evaluation constraints to avoid post-hoc metric selection bias.

\textbf{Depth guidance (advanced).} At the advanced level, the same concept should be treated as a research variable: students should formulate competing hypotheses, design controlled ablations, and reason about how uncertainty in one subsystem propagates into frontier movement.

\textbf{Operational checklist.} Before accepting any run result for this module, verify: (i) metric definitions are unchanged across runs, (ii) scenario and seed lists are fixed, (iii) checkpoint metadata is parseable by automated evaluators, (iv) run logs include safety channels, and (v) threshold tier classification is reported together with raw reward and cost.

\begin{itemize}[leftmargin=*]
  \item Principle statement: unit and smoke tests protect refactors.
  \item Primary challenge: unvalidated merges can silently break training.
  \item Mitigation path: test env determinism, update-step sanity, and CLI smoke paths.
  \item Verification channels: reward, cost, lambda, risk, block rate.
  \item Success criterion: positive frontier trend with stable safety traces.
  \item Documentation rule: report both mean and seed-level outcomes.
\end{itemize}


\subsubsection{Module 47: Statistical Reporting (Foundational View)}
\textbf{Principle.} means alone are insufficient for high-variance domains. In AMASA-Complex v2 this appears as a constraint on both algorithm design and experiment process. The key point is that no isolated formula is sufficient; correctness emerges from interaction between objective, data regime, guard behavior, and reproducibility checks.

\textbf{Failure mode.} A common practical breakdown is single-seed claims are unreliable. This usually manifests as deceptive short-run gains followed by cross-seed instability. In safety-critical tasks, this instability must be interpreted as a primary failure signal even when average reward appears competitive.

\textbf{Mitigation.} The project mitigation is to report seed-wise spread, scenario means, and frontier ranking. This mitigation is encoded in configuration and code so that it can be repeated consistently across seeds and scenarios. It is also paired with evaluation constraints to avoid post-hoc metric selection bias.

\textbf{Depth guidance (foundational).} At the foundational level, students should be able to restate this concept formally, map it to concrete code paths, and identify the minimum telemetry needed to validate implementation correctness.

\textbf{Operational checklist.} Before accepting any run result for this module, verify: (i) metric definitions are unchanged across runs, (ii) scenario and seed lists are fixed, (iii) checkpoint metadata is parseable by automated evaluators, (iv) run logs include safety channels, and (v) threshold tier classification is reported together with raw reward and cost.

\begin{itemize}[leftmargin=*]
  \item Principle statement: means alone are insufficient for high-variance domains.
  \item Primary challenge: single-seed claims are unreliable.
  \item Mitigation path: report seed-wise spread, scenario means, and frontier ranking.
  \item Verification channels: reward, cost, lambda, risk, block rate.
  \item Success criterion: positive frontier trend with stable safety traces.
  \item Documentation rule: report both mean and seed-level outcomes.
\end{itemize}


\subsubsection{Module 48: Statistical Reporting (Advanced Research View)}
\textbf{Principle.} means alone are insufficient for high-variance domains. In AMASA-Complex v2 this appears as a constraint on both algorithm design and experiment process. The key point is that no isolated formula is sufficient; correctness emerges from interaction between objective, data regime, guard behavior, and reproducibility checks.

\textbf{Failure mode.} A common practical breakdown is single-seed claims are unreliable. This usually manifests as deceptive short-run gains followed by cross-seed instability. In safety-critical tasks, this instability must be interpreted as a primary failure signal even when average reward appears competitive.

\textbf{Mitigation.} The project mitigation is to report seed-wise spread, scenario means, and frontier ranking. This mitigation is encoded in configuration and code so that it can be repeated consistently across seeds and scenarios. It is also paired with evaluation constraints to avoid post-hoc metric selection bias.

\textbf{Depth guidance (advanced).} At the advanced level, the same concept should be treated as a research variable: students should formulate competing hypotheses, design controlled ablations, and reason about how uncertainty in one subsystem propagates into frontier movement.

\textbf{Operational checklist.} Before accepting any run result for this module, verify: (i) metric definitions are unchanged across runs, (ii) scenario and seed lists are fixed, (iii) checkpoint metadata is parseable by automated evaluators, (iv) run logs include safety channels, and (v) threshold tier classification is reported together with raw reward and cost.

\begin{itemize}[leftmargin=*]
  \item Principle statement: means alone are insufficient for high-variance domains.
  \item Primary challenge: single-seed claims are unreliable.
  \item Mitigation path: report seed-wise spread, scenario means, and frontier ranking.
  \item Verification channels: reward, cost, lambda, risk, block rate.
  \item Success criterion: positive frontier trend with stable safety traces.
  \item Documentation rule: report both mean and seed-level outcomes.
\end{itemize}


\subsubsection{Module 49: Human-in-the-Loop Safety (Foundational View)}
\textbf{Principle.} operators need transparent intervention rationale. In AMASA-Complex v2 this appears as a constraint on both algorithm design and experiment process. The key point is that no isolated formula is sufficient; correctness emerges from interaction between objective, data regime, guard behavior, and reproducibility checks.

\textbf{Failure mode.} A common practical breakdown is opaque blocks reduce trust during deployment. This usually manifests as deceptive short-run gains followed by cross-seed instability. In safety-critical tasks, this instability must be interpreted as a primary failure signal even when average reward appears competitive.

\textbf{Mitigation.} The project mitigation is to attach simple natural-language explanations to shield decisions. This mitigation is encoded in configuration and code so that it can be repeated consistently across seeds and scenarios. It is also paired with evaluation constraints to avoid post-hoc metric selection bias.

\textbf{Depth guidance (foundational).} At the foundational level, students should be able to restate this concept formally, map it to concrete code paths, and identify the minimum telemetry needed to validate implementation correctness.

\textbf{Operational checklist.} Before accepting any run result for this module, verify: (i) metric definitions are unchanged across runs, (ii) scenario and seed lists are fixed, (iii) checkpoint metadata is parseable by automated evaluators, (iv) run logs include safety channels, and (v) threshold tier classification is reported together with raw reward and cost.

\begin{itemize}[leftmargin=*]
  \item Principle statement: operators need transparent intervention rationale.
  \item Primary challenge: opaque blocks reduce trust during deployment.
  \item Mitigation path: attach simple natural-language explanations to shield decisions.
  \item Verification channels: reward, cost, lambda, risk, block rate.
  \item Success criterion: positive frontier trend with stable safety traces.
  \item Documentation rule: report both mean and seed-level outcomes.
\end{itemize}


\subsubsection{Module 50: Human-in-the-Loop Safety (Advanced Research View)}
\textbf{Principle.} operators need transparent intervention rationale. In AMASA-Complex v2 this appears as a constraint on both algorithm design and experiment process. The key point is that no isolated formula is sufficient; correctness emerges from interaction between objective, data regime, guard behavior, and reproducibility checks.

\textbf{Failure mode.} A common practical breakdown is opaque blocks reduce trust during deployment. This usually manifests as deceptive short-run gains followed by cross-seed instability. In safety-critical tasks, this instability must be interpreted as a primary failure signal even when average reward appears competitive.

\textbf{Mitigation.} The project mitigation is to attach simple natural-language explanations to shield decisions. This mitigation is encoded in configuration and code so that it can be repeated consistently across seeds and scenarios. It is also paired with evaluation constraints to avoid post-hoc metric selection bias.

\textbf{Depth guidance (advanced).} At the advanced level, the same concept should be treated as a research variable: students should formulate competing hypotheses, design controlled ablations, and reason about how uncertainty in one subsystem propagates into frontier movement.

\textbf{Operational checklist.} Before accepting any run result for this module, verify: (i) metric definitions are unchanged across runs, (ii) scenario and seed lists are fixed, (iii) checkpoint metadata is parseable by automated evaluators, (iv) run logs include safety channels, and (v) threshold tier classification is reported together with raw reward and cost.

\begin{itemize}[leftmargin=*]
  \item Principle statement: operators need transparent intervention rationale.
  \item Primary challenge: opaque blocks reduce trust during deployment.
  \item Mitigation path: attach simple natural-language explanations to shield decisions.
  \item Verification channels: reward, cost, lambda, risk, block rate.
  \item Success criterion: positive frontier trend with stable safety traces.
  \item Documentation rule: report both mean and seed-level outcomes.
\end{itemize}


\subsubsection{Module 51: Dataset Curation for Offline RL (Foundational View)}
\textbf{Principle.} data diversity is as important as data volume. In AMASA-Complex v2 this appears as a constraint on both algorithm design and experiment process. The key point is that no isolated formula is sufficient; correctness emerges from interaction between objective, data regime, guard behavior, and reproducibility checks.

\textbf{Failure mode.} A common practical breakdown is narrow behavior coverage limits policy headroom. This usually manifests as deceptive short-run gains followed by cross-seed instability. In safety-critical tasks, this instability must be interpreted as a primary failure signal even when average reward appears competitive.

\textbf{Mitigation.} The project mitigation is to include near-boundary trajectories and controlled failures. This mitigation is encoded in configuration and code so that it can be repeated consistently across seeds and scenarios. It is also paired with evaluation constraints to avoid post-hoc metric selection bias.

\textbf{Depth guidance (foundational).} At the foundational level, students should be able to restate this concept formally, map it to concrete code paths, and identify the minimum telemetry needed to validate implementation correctness.

\textbf{Operational checklist.} Before accepting any run result for this module, verify: (i) metric definitions are unchanged across runs, (ii) scenario and seed lists are fixed, (iii) checkpoint metadata is parseable by automated evaluators, (iv) run logs include safety channels, and (v) threshold tier classification is reported together with raw reward and cost.

\begin{itemize}[leftmargin=*]
  \item Principle statement: data diversity is as important as data volume.
  \item Primary challenge: narrow behavior coverage limits policy headroom.
  \item Mitigation path: include near-boundary trajectories and controlled failures.
  \item Verification channels: reward, cost, lambda, risk, block rate.
  \item Success criterion: positive frontier trend with stable safety traces.
  \item Documentation rule: report both mean and seed-level outcomes.
\end{itemize}


\subsubsection{Module 52: Dataset Curation for Offline RL (Advanced Research View)}
\textbf{Principle.} data diversity is as important as data volume. In AMASA-Complex v2 this appears as a constraint on both algorithm design and experiment process. The key point is that no isolated formula is sufficient; correctness emerges from interaction between objective, data regime, guard behavior, and reproducibility checks.

\textbf{Failure mode.} A common practical breakdown is narrow behavior coverage limits policy headroom. This usually manifests as deceptive short-run gains followed by cross-seed instability. In safety-critical tasks, this instability must be interpreted as a primary failure signal even when average reward appears competitive.

\textbf{Mitigation.} The project mitigation is to include near-boundary trajectories and controlled failures. This mitigation is encoded in configuration and code so that it can be repeated consistently across seeds and scenarios. It is also paired with evaluation constraints to avoid post-hoc metric selection bias.

\textbf{Depth guidance (advanced).} At the advanced level, the same concept should be treated as a research variable: students should formulate competing hypotheses, design controlled ablations, and reason about how uncertainty in one subsystem propagates into frontier movement.

\textbf{Operational checklist.} Before accepting any run result for this module, verify: (i) metric definitions are unchanged across runs, (ii) scenario and seed lists are fixed, (iii) checkpoint metadata is parseable by automated evaluators, (iv) run logs include safety channels, and (v) threshold tier classification is reported together with raw reward and cost.

\begin{itemize}[leftmargin=*]
  \item Principle statement: data diversity is as important as data volume.
  \item Primary challenge: narrow behavior coverage limits policy headroom.
  \item Mitigation path: include near-boundary trajectories and controlled failures.
  \item Verification channels: reward, cost, lambda, risk, block rate.
  \item Success criterion: positive frontier trend with stable safety traces.
  \item Documentation rule: report both mean and seed-level outcomes.
\end{itemize}


\subsubsection{Module 53: Scaling to Real Robotics (Foundational View)}
\textbf{Principle.} sim-to-real requires robust uncertainty handling. In AMASA-Complex v2 this appears as a constraint on both algorithm design and experiment process. The key point is that no isolated formula is sufficient; correctness emerges from interaction between objective, data regime, guard behavior, and reproducibility checks.

\textbf{Failure mode.} A common practical breakdown is sim-only policies exploit unrealistic shortcuts. This usually manifests as deceptive short-run gains followed by cross-seed instability. In safety-critical tasks, this instability must be interpreted as a primary failure signal even when average reward appears competitive.

\textbf{Mitigation.} The project mitigation is to combine randomization, conservative objectives, and explicit safety filters. This mitigation is encoded in configuration and code so that it can be repeated consistently across seeds and scenarios. It is also paired with evaluation constraints to avoid post-hoc metric selection bias.

\textbf{Depth guidance (foundational).} At the foundational level, students should be able to restate this concept formally, map it to concrete code paths, and identify the minimum telemetry needed to validate implementation correctness.

\textbf{Operational checklist.} Before accepting any run result for this module, verify: (i) metric definitions are unchanged across runs, (ii) scenario and seed lists are fixed, (iii) checkpoint metadata is parseable by automated evaluators, (iv) run logs include safety channels, and (v) threshold tier classification is reported together with raw reward and cost.

\begin{itemize}[leftmargin=*]
  \item Principle statement: sim-to-real requires robust uncertainty handling.
  \item Primary challenge: sim-only policies exploit unrealistic shortcuts.
  \item Mitigation path: combine randomization, conservative objectives, and explicit safety filters.
  \item Verification channels: reward, cost, lambda, risk, block rate.
  \item Success criterion: positive frontier trend with stable safety traces.
  \item Documentation rule: report both mean and seed-level outcomes.
\end{itemize}


\subsubsection{Module 54: Scaling to Real Robotics (Advanced Research View)}
\textbf{Principle.} sim-to-real requires robust uncertainty handling. In AMASA-Complex v2 this appears as a constraint on both algorithm design and experiment process. The key point is that no isolated formula is sufficient; correctness emerges from interaction between objective, data regime, guard behavior, and reproducibility checks.

\textbf{Failure mode.} A common practical breakdown is sim-only policies exploit unrealistic shortcuts. This usually manifests as deceptive short-run gains followed by cross-seed instability. In safety-critical tasks, this instability must be interpreted as a primary failure signal even when average reward appears competitive.

\textbf{Mitigation.} The project mitigation is to combine randomization, conservative objectives, and explicit safety filters. This mitigation is encoded in configuration and code so that it can be repeated consistently across seeds and scenarios. It is also paired with evaluation constraints to avoid post-hoc metric selection bias.

\textbf{Depth guidance (advanced).} At the advanced level, the same concept should be treated as a research variable: students should formulate competing hypotheses, design controlled ablations, and reason about how uncertainty in one subsystem propagates into frontier movement.

\textbf{Operational checklist.} Before accepting any run result for this module, verify: (i) metric definitions are unchanged across runs, (ii) scenario and seed lists are fixed, (iii) checkpoint metadata is parseable by automated evaluators, (iv) run logs include safety channels, and (v) threshold tier classification is reported together with raw reward and cost.

\begin{itemize}[leftmargin=*]
  \item Principle statement: sim-to-real requires robust uncertainty handling.
  \item Primary challenge: sim-only policies exploit unrealistic shortcuts.
  \item Mitigation path: combine randomization, conservative objectives, and explicit safety filters.
  \item Verification channels: reward, cost, lambda, risk, block rate.
  \item Success criterion: positive frontier trend with stable safety traces.
  \item Documentation rule: report both mean and seed-level outcomes.
\end{itemize}


\subsubsection{Module 55: Course Integration Strategy (Foundational View)}
\textbf{Principle.} students benefit from end-to-end systems, not isolated scripts. In AMASA-Complex v2 this appears as a constraint on both algorithm design and experiment process. The key point is that no isolated formula is sufficient; correctness emerges from interaction between objective, data regime, guard behavior, and reproducibility checks.

\textbf{Failure mode.} A common practical breakdown is fragmented assignments miss cross-stage coupling. This usually manifests as deceptive short-run gains followed by cross-seed instability. In safety-critical tasks, this instability must be interpreted as a primary failure signal even when average reward appears competitive.

\textbf{Mitigation.} The project mitigation is to evaluate architecture, safety, and analysis quality jointly. This mitigation is encoded in configuration and code so that it can be repeated consistently across seeds and scenarios. It is also paired with evaluation constraints to avoid post-hoc metric selection bias.

\textbf{Depth guidance (foundational).} At the foundational level, students should be able to restate this concept formally, map it to concrete code paths, and identify the minimum telemetry needed to validate implementation correctness.

\textbf{Operational checklist.} Before accepting any run result for this module, verify: (i) metric definitions are unchanged across runs, (ii) scenario and seed lists are fixed, (iii) checkpoint metadata is parseable by automated evaluators, (iv) run logs include safety channels, and (v) threshold tier classification is reported together with raw reward and cost.

\begin{itemize}[leftmargin=*]
  \item Principle statement: students benefit from end-to-end systems, not isolated scripts.
  \item Primary challenge: fragmented assignments miss cross-stage coupling.
  \item Mitigation path: evaluate architecture, safety, and analysis quality jointly.
  \item Verification channels: reward, cost, lambda, risk, block rate.
  \item Success criterion: positive frontier trend with stable safety traces.
  \item Documentation rule: report both mean and seed-level outcomes.
\end{itemize}


\subsubsection{Module 56: Course Integration Strategy (Advanced Research View)}
\textbf{Principle.} students benefit from end-to-end systems, not isolated scripts. In AMASA-Complex v2 this appears as a constraint on both algorithm design and experiment process. The key point is that no isolated formula is sufficient; correctness emerges from interaction between objective, data regime, guard behavior, and reproducibility checks.

\textbf{Failure mode.} A common practical breakdown is fragmented assignments miss cross-stage coupling. This usually manifests as deceptive short-run gains followed by cross-seed instability. In safety-critical tasks, this instability must be interpreted as a primary failure signal even when average reward appears competitive.

\textbf{Mitigation.} The project mitigation is to evaluate architecture, safety, and analysis quality jointly. This mitigation is encoded in configuration and code so that it can be repeated consistently across seeds and scenarios. It is also paired with evaluation constraints to avoid post-hoc metric selection bias.

\textbf{Depth guidance (advanced).} At the advanced level, the same concept should be treated as a research variable: students should formulate competing hypotheses, design controlled ablations, and reason about how uncertainty in one subsystem propagates into frontier movement.

\textbf{Operational checklist.} Before accepting any run result for this module, verify: (i) metric definitions are unchanged across runs, (ii) scenario and seed lists are fixed, (iii) checkpoint metadata is parseable by automated evaluators, (iv) run logs include safety channels, and (v) threshold tier classification is reported together with raw reward and cost.

\begin{itemize}[leftmargin=*]
  \item Principle statement: students benefit from end-to-end systems, not isolated scripts.
  \item Primary challenge: fragmented assignments miss cross-stage coupling.
  \item Mitigation path: evaluate architecture, safety, and analysis quality jointly.
  \item Verification channels: reward, cost, lambda, risk, block rate.
  \item Success criterion: positive frontier trend with stable safety traces.
  \item Documentation rule: report both mean and seed-level outcomes.
\end{itemize}

\section{Empirical Results from Local Runs}
\subsection{Smoke Benchmark Matrix}
The smoke benchmark matrix completed 30 jobs and exported aggregate summaries. Means by algorithm and scenario are shown in Table~\ref{tab:benchmeans}.

\begin{table}[!t]
\centering
\caption{Benchmark means from \texttt{results/bench\_smoke/benchmark\_summary.csv}}
\label{tab:benchmeans}
\begin{tabular}{llrr}
\toprule
Algo & Scenario & Avg Reward & Avg Cost \\
\midrule
CQL & Nominal & 594.27 & 0.096 \\
CQL & Perturbed & 440.04 & 0.118 \\
IQL & Nominal & 439.77 & 0.057 \\
IQL & Perturbed & 746.91 & 0.099 \\
\midrule
SAC-Lag & Nominal & -227.40 & 0.908 \\
SAC-Lag & Perturbed & 50.96 & 0.730 \\
SAC-Lag & Adversarial & 236.96 & 0.616 \\
\midrule
PPO-Lag & Nominal & -218.26 & 0.788 \\
PPO-Lag & Perturbed & -203.95 & 0.918 \\
PPO-Lag & Adversarial & -94.16 & 0.797 \\
\bottomrule
\end{tabular}
\end{table}

\begin{figure}[!t]
\centering
\includegraphics[width=0.98\columnwidth]{plots/global_pareto.png}
\caption{Global reward-cost Pareto across smoke benchmark jobs.}
\end{figure}

\begin{figure}[!t]
\centering
\subfloat[Nominal]{\includegraphics[width=0.31\columnwidth]{plots/pareto_nominal.png}}
\hfill
\subfloat[Perturbed]{\includegraphics[width=0.31\columnwidth]{plots/pareto_perturbed.png}}
\hfill
\subfloat[Adversarial]{\includegraphics[width=0.31\columnwidth]{plots/pareto_adversarial.png}}
\caption{Scenario-level reward-cost curves.}
\end{figure}

\subsection{Expanded PID Sweep}
The PID grid was expanded to $k_p\in\{0.4,0.5,0.55,0.6,0.7\}$ and $k_d\in\{0.05,0.08,0.10,0.12\}$ with $k_i=0.08$ (20 runs total).

\begin{table}[!t]
\centering
\caption{Top PID settings by frontier score ($R-300C$)}
\begin{tabular}{rrrrr}
\toprule
$k_p$ & $k_d$ & Reward & Cost & Frontier \\
\midrule
0.60 & 0.08 & 419.65 & 0.685 & 214.15 \\
0.55 & 0.10 & 412.70 & 0.698 & 203.20 \\
0.70 & 0.10 & 363.79 & 0.848 & 109.29 \\
0.60 & 0.05 & 332.65 & 0.770 & 101.65 \\
0.50 & 0.10 & -99.11 & 0.778 & -332.61 \\
0.40 & 0.12 & -131.82 & 0.835 & -382.32 \\
0.60 & 0.12 & -131.54 & 0.852 & -387.04 \\
\bottomrule
\end{tabular}
\end{table}

\begin{figure}[!t]
\centering
\subfloat[Reward]{\includegraphics[width=0.31\columnwidth]{plots/pid_heatmap_reward.png}}
\hfill
\subfloat[Cost]{\includegraphics[width=0.31\columnwidth]{plots/pid_heatmap_cost.png}}
\hfill
\subfloat[Frontier]{\includegraphics[width=0.31\columnwidth]{plots/pid_heatmap_frontier.png}}
\caption{Expanded PID heatmaps.}
\end{figure}

\subsection{Safety Timeline}
\begin{figure}[!t]
\centering
\includegraphics[width=0.98\columnwidth]{plots/safety_timeline_best.png}
\caption{Best-frontier safety timeline showing lambda, risk score, and block-rate trend.}
\end{figure}

\subsection{Threshold Tier Interpretation}
Threshold counts from benchmark evaluation are:
\begin{itemize}[leftmargin=*]
  \item stretch: 8
  \item target: 3
  \item minimum: 2
  \item below\_min: 17
  \item offline\_gate\_pass: 7
  \item offline\_gate\_fail: 5
\end{itemize}
These counts show that strong offline starts are achievable in smoke budgets, while online safety-constrained optimization still requires longer adaptation windows and better dual dynamics calibration.

\section{Engineering Lessons and Full-Preset Roadmap}
\subsection{What Worked}
\begin{itemize}[leftmargin=*]
  \item Config-first orchestration and schema checks reduced silent experiment drift.
  \item Resumable benchmark/sweep execution made long runs tractable.
  \item Safety timeline exports improved diagnosis of guard saturation behavior.
  \item Expanded PID grid identified better frontier regions than the initial small grid.
\end{itemize}

\subsection{What Still Limits Performance}
\begin{itemize}[leftmargin=*]
  \item Online methods can remain trapped in high-lambda conservative regimes.
  \item Cost critics can be under-calibrated in early online phases.
  \item Fixed cost limits may be too strict during initial adaptation.
  \item On-policy constrained updates remain sensitive to advantage variance.
\end{itemize}

\subsection{Recommended Full-Preset Study Plan}
\begin{enumerate}[leftmargin=*]
  \item Increase offline warm-start duration for CQL and IQL.
  \item Use best PID seeds from sweep as initialization for long runs.
  \item Run staged cost-limit schedules in early online training.
  \item Execute ablation matrix over guard components.
  \item Report tier movement and frontier movement jointly.
\end{enumerate}

\section{Conclusion}
AMASA-Complex v2 now provides a complete, deeply analyzable project for safe RL in a high-stakes domain. The system is not only algorithmically richer but also operationally stronger: configuration discipline, compatibility guarantees, robust logging, reproducible matrix execution, and report-ready visual analytics are integrated. The project therefore meets both educational and engineering objectives and provides a reliable foundation for longer full-preset studies.

\appendices
\section{Command Cookbook}
\begin{lstlisting}[language=bash,basicstyle=\ttfamily\footnotesize,breaklines=true]
PYTHONPATH=/Users/tahamajs/Documents/uni/DRL \
pytest -q /Users/tahamajs/Documents/uni/DRL/projects/amasa_clean/tests

PYTHONPATH=/Users/tahamajs/Documents/uni/DRL \
python3.10 -m projects.amasa_clean.scripts.run_experiment \
  --mode benchmark \
  --config projects/amasa_clean/configs/base.yaml \
  --preset smoke \
  --out_dir projects/amasa_clean/results/bench_smoke \
  --dataset data/amasa_offline.npz

PYTHONPATH=/Users/tahamajs/Documents/uni/DRL \
python3.10 -m projects.amasa_clean.scripts.run_experiment \
  --mode pid_sweep \
  --config projects/amasa_clean/configs/base.yaml \
  --preset smoke \
  --algo sac_lag \
  --scenario perturbed \
  --out_dir projects/amasa_clean/results/pid_smoke \
  --dataset data/amasa_offline.npz
\end{lstlisting}

\section{Overall Means by Algorithm}
\begin{table}[!h]
\centering
\caption{Overall mean reward and cost in smoke benchmark}
\begin{tabular}{lrr}
\toprule
Algorithm & Mean Reward & Mean Cost \\
\midrule
cql & 517.16 & 0.107 \\
iql & 593.34 & 0.078 \\
ppo\_lag & -172.12 & 0.834 \\
sac\_lag & 20.17 & 0.751 \\
\bottomrule
\end{tabular}
\end{table}

\section{Verification Checklist}
\begin{itemize}[leftmargin=*]
  \item Unit and smoke tests pass.
  \item Benchmark summary and threshold CSVs exist.
  \item PID sweep summary and threshold CSVs exist.
  \item Pareto, per-scenario, PID heatmap, and safety timeline figures exist.
  \item Report sources and PDFs are generated from latest outputs.
\end{itemize}

\end{document}
